\documentclass[11pt]{article}
\usepackage[margin=1in]{geometry}
\usepackage{amsmath,amssymb,amsthm,mathtools,mathrsfs,bm}
\usepackage{enumitem}
\usepackage{booktabs}
\usepackage{hyperref}
\usepackage{microtype}

\allowdisplaybreaks

\newtheorem{theorem}{Theorem}[section]
\newtheorem{proposition}[theorem]{Proposition}
\newtheorem{lemma}[theorem]{Lemma}
\newtheorem{corollary}[theorem]{Corollary}
\theoremstyle{definition}
\newtheorem{definition}[theorem]{Definition}
\newtheorem{assumption}[theorem]{Assumption}
\theoremstyle{remark}
\newtheorem{remark}[theorem]{Remark}

\newcommand{\R}{\mathbb R}
\newcommand{\C}{\mathbb C}
\newcommand{\N}{\mathbb N}
\newcommand{\Z}{\mathbb Z}
\newcommand{\Tset}{\{0,1,\dots,T\}}
\newcommand{\Law}{\mathrm{Law}}
\newcommand{\E}{\mathbb E}
\newcommand{\KL}{\mathrm{KL}}
\newcommand{\Ent}{\mathrm{Ent}}
\newcommand{\TV}{\mathrm{TV}}
\newcommand{\mcN}{\mathcal N}
\newcommand{\mcE}{\mathcal E}
\newcommand{\mcL}{\mathcal L}
\newcommand{\mcQ}{\mathcal Q}
\newcommand{\mcC}{\mathcal C}
\newcommand{\mcP}{\mathcal P}
\newcommand{\mcK}{\mathcal K}
\newcommand{\mcV}{\mathcal V}
\newcommand{\mcX}{\mathcal X}
\newcommand{\mcB}{\mathcal B}
\newcommand{\mcR}{\mathcal R}
\newcommand{\dd}{\,\mathrm d}
\newcommand{\tr}{\operatorname{tr}}
\newcommand{\supp}{\operatorname{supp}}
\newcommand{\one}{\mathbf 1}
\newcommand{\fraksu}{\mathfrak{su}(3)}
\newcommand{\mathfrakh}{\mathfrak h_3}

\title{Finite Yang--Mills Sector on the Realized Fractal Regulator\\
with a Classical Continuum Comparison}
\author{}
\date{}

\begin{document}
\maketitle

\begin{abstract}
This article rewrites the Yang--Mills sector of the particle system furnished
by \texttt{01\_convergence} in Fractal-Gas form. The foundational probability
space is the swarm path law itself. From each realized recorded run we extract
the Fractal regulator, namely a random oriented \(2\)-complex with CST, IG,
and IA edges and interaction plaquettes. Gauge variables are assigned directly
from realized local state data on that random complex, so plaquette holonomies,
Wilson loops, and the Wilson action become genuine random observables. The main
finite theorem defines the induced gauge law as the pushforward of the swarm
path law through the regulator-and-gauge construction map. Under a local
refinement and small-field regime, expectations of localized Wilson observables
recover the continuum Yang--Mills density on compact windows.
\end{abstract}

\tableofcontents

\section{Introduction}

The finite regulator in this paper is the realized Fractal Set generated by the
stochastic swarm itself, and the finite gauge sector is defined directly on
that random geometry.

For each finite swarm size \(N\) and observation horizon \(T\), the recorded
run history determines a random oriented \(2\)-complex
\[
\mathfrak F^{(N,T)}(\omega)
=
\bigl(\mcV^{(N,T)}(\omega),\mcE_{\mathrm{CST}}^{(N,T)}(\omega),
\mcE_{\mathrm{IG}}^{(N,T)}(\omega),\mcE_{\mathrm{IA}}^{(N,T)}(\omega),
\mcP_{\mathrm{FG}}^{(N,T)}(\omega)\bigr),
\]
whose nodes are recorded pre/clone/final events, whose edge families encode
temporal evolution, companion couplings, and influence attribution, and whose
plaquettes are the interaction \(4\)-cycles obtained from pairs of opposite
interaction triangles. On that realized regulator we build group-valued edge
transports \(U_e(\omega)\in SU(3)\), plaquette holonomies
\[
U_P(\omega)=\prod_{e\in\partial P}U_e(\omega),
\]
Wilson loop observables
\[
W_\gamma(\omega):=\frac13\tr\!\Bigl(\prod_{e\in\gamma}U_e(\omega)\Bigr),
\]
and the random Wilson action
\[
S_{\mathrm W}^{(N,T)}(\omega)
:=
\frac{6}{g_3^2}\sum_{P\in\mcP_{\mathrm{FG}}^{(N,T)}(\omega)}
\left(1-\frac13\Re\tr U_P(\omega)\right).
\]

The paper is organized around four steps:
\begin{enumerate}[leftmargin=1.5em]
\item import from \texttt{01\_convergence} the stochastic substrate and the
stationary mean-field field;
\item extract from each realized swarm history the Fractal regulator and assign
realized edge transports from local state data;
\item define the induced finite gauge law as the pushforward of the swarm path
law through the regulator-and-gauge construction map;
\item identify the local continuum density governing the expected small-loop
limit of localized Wilson observables.
\end{enumerate}

\section{Particle-system preliminaries}

\begin{definition}[Trajectory notation for bridge quantities]
\label{def:bridge-notation}
For each family index $N$, and each observation time
$t\in\{0,\dots,T\}$, define
\[
Z^{(N)}(t):=(X_1^{(N)}(t),V_1^{(N)}(t),\dots,X_N^{(N)}(t),V_N^{(N)}(t))
\in(\R^8)^N,
\]
and
\[
\mu_t^N:=\frac1N\sum_{i=1}^N\delta_{(X_i^{(N)}(t),V_i^{(N)}(t))},
\qquad
\rho_t^N:=\Law(Z^{(N)}(t)),
\qquad
\rho_{t,1}^N:=\text{first marginal of }\rho_t^N.
\]
Let \((\bar X_i,\bar V_i)\) be an exchangeable nonlinear mean-field family with
common one-particle law \(f_t\), and empirical average
\[
\bar\mu_t^N:=\frac1N\sum_{i=1}^N\delta_{(\bar X_i(t),\bar V_i(t))}.
\]
Finite invariant and mean-field stationary laws are denoted
\[
\pi_N\in\mcP((\R^8)^N),\quad
\pi_{N,1}\in\mcP(\R^8),\quad
\pi_\infty\in\mcP(\R^8).
\]
\end{definition}

\begin{assumption}[Kinetic-core input translated from 01]
\label{ass:01-core}
\begin{enumerate}[leftmargin=1.5em,label=\textnormal{(\roman*)}]
\item\textbf{Dimension and confining core.}
The position and velocity variables are in $d=4$.
The confining potential \(U\in C^2(\R^4)\) is globally regular and strongly
coercive:
\begin{align*}
\|\nabla U(x)-\nabla U(y)\|&\le L_U\|x-y\|,\\
\langle\nabla U(x)-\nabla U(y),x-y\rangle&\ge L_*\|x-y\|^2,\\
\langle x,\nabla U(x)\rangle&\ge m_U\|x\|^2-\beta_U.
\end{align*}
\item\textbf{Regularized viscous kernel and companion rule.}
For each \(N\) the interaction scale \(\rho_N>0\) defines
\[
K_{\rho_N}(x,y):=\kappa_0+\exp\!\left(-\frac{\|x-y\|^2}{2\rho_N^2}\right),
\qquad
K_{\rho_N}(x,y)\ge\kappa_0>0.
\]
\item\textbf{Core drift and noise scales.}
The finite-particle kinetic backbone has damping \(\gamma>0\), noise strength
\(\sigma>0\), and viscous coupling \(\nu_{\mathrm{visc}}\ge0\), with Brownian forcing and dissipative drift as in the kinetic core of \texttt{01\_convergence}.
\item\textbf{Velocity observation.}
The velocity-squashing map is
\[
\psi(v):=\frac{V_{\mathrm{alg}}}{V_{\mathrm{alg}}+\|v\|}\,v,\qquad V_{\mathrm{alg}}>0.
\]
\end{enumerate}
\end{assumption}

\paragraph{Bridge assumptions}
\begin{assumption}[Particle-to-mean-field transfer and compactness]
\label{ass:bridge-backbone}
Let $\{\mcQ^{(N)}\}_{N\ge N_0}$ be the sequence of trajectory systems from this
paper, with empirical laws $\rho_t^N$ on $(\R^4\times\R^4)^N$ and mean-field
law $f_t$ on $\R^4\times\R^4$.
The following constants $L_U,L_*,\gamma,\sigma,\nu_{\mathrm{visc}},C_T,C_{\mathrm{LSI}},
C_{\mathrm{tail}}>0$ and random variables are fixed from the outset.
\begin{enumerate}[leftmargin=1.5em,label=\textnormal{(\roman*)}]
\item \textbf{Confinement and kernel structure.}
The confining potential and interaction kernel obey strong regularity and confining
bounds:
\[
U\in C^2,\qquad \|\nabla U(x)-\nabla U(y)\|\le L_U\|x-y\|,\qquad
\langle \nabla U(x)-\nabla U(y),x-y\rangle\ge L_*\|x-y\|^2,
\]
and
\[
K_{\rho}(x,y)=\kappa_0+\exp\!\left(-\frac{\|x-y\|^2}{2\rho^2}\right),
\qquad
K_{\rho}(x,y)\ge \kappa_0>0.
\]
\item \textbf{Finite-time concentration / chaos.}
For every finite horizon $T>0$ and every bounded-Lipschitz observable
$\Phi$, there is $C_T<\infty$ such that
\[
\KL(\rho_t^N\|f_t^{\otimes N})\le C_T,\qquad
d_{\TV}(\rho_{t,1}^N,f_t)\le\frac{C_T}{\sqrt N},
\qquad 0\le t\le T,
\]
and the first-particle Wasserstein transport estimate
\[
\sup_{t\in[0,T]}\E\|X_1^{(N)}(t)-\bar X_1(t)\|^2+\E\|V_1^{(N)}(t)-\bar V_1(t)\|^2
\le \frac{C_T}{N}
\]
holds for the finite/mean-field coupling.
\item \textbf{Stationary thermodynamic limit and moment control.}
\[
\pi_{N,1}\Rightarrow\pi_\infty,
\qquad
\sup_N\int(\|x\|^2+\|v\|^2)\,\pi_{N,1}(dx\,dv)<\infty.
\]
\item \textbf{Direct mean-field functional inequality and wall-to-infinity tail control.}
The stationary mean-field law satisfies
\[
\Ent_{\pi_\infty}(g^2)\le C_{\mathrm{LSI}}
\int\Gamma_{\mathrm{hypo}}^\infty(g,g)\,d\pi_\infty
\]
for smooth compactly supported \(g\). Moreover, for every bounded continuous
observable $\Phi:\R^8\to\R$ and every recentering sequence $a_R\in\R^4$ with
$\|a_R\|\to\infty$,
\[
\lim_{R\to\infty}\lim_{N\to\infty}
\int\Phi(x-a_R,v)\,\pi_{N,1}^{a_R}(dx\,dv)
=
\int\Phi(x,v)\,\pi_{\infty}^0(dx\,dv),
\]
where $\pi_{N,1}^{a_R}$ and $\pi_\infty^{a_R}$ are finite-$N$ and mean-field stationary
laws under translated confining envelopes.
\end{enumerate}
\end{assumption}

\begin{theorem}[Bridge package from 01\_convergence]
\label{thm:bridge-from-01}
Assumption~\ref{ass:01-core} yields all four items of
Assumption~\ref{ass:bridge-backbone}, with explicit estimates as follows:
\begin{enumerate}[leftmargin=1.5em,label=\textnormal{(\roman*)}]
\item Finite-time chaos and finite-time coupling:
for every finite horizon $T>0$ there is $C_T<\infty$ such that
\begin{align*}
\KL(\rho_t^N\|f_t^{\otimes N})&\leq C_T,\\
d_{\TV}(\rho_{t,1}^N,f_t)&\leq \frac{C_T}{\sqrt N},\\
\sup_{t\in[0,T]}
\E\!\left(
\|X_1^{(N)}(t)-\bar X_1(t)\|^2
+\|V_1^{(N)}(t)-\bar V_1(t)\|^2
\right)
&\leq \frac{C_T}{N}.
\end{align*}
\item Direct mean-field regularization:
the stationary mean-field law is unique, smooth, strictly positive, and satisfies
\[
\Ent_{\pi_\infty}(g^2)
\le
C_{\mathrm{LSI}}\int_{\R^8}\Gamma_{\mathrm{hypo}}^\infty(g,g)\,d\pi_\infty.
\]
\item Stationary compactness and thermodynamic limit:
there are constants $C_2<\infty$, $N_0<\infty$ such that
\[
\sup_{N\ge N_0}\int_{\R^8}(\|x\|^2+\|v\|^2)\,\pi_{N,1}(dx\,dv)\le C_2,
\]
and
\[
\pi_{N,1}\Rightarrow\pi_\infty
\]
in $\mcP(\R^8)$ as $N\to\infty$.
\item Recentered translated limit:
for every bounded $\Phi:\R^8\to\R$ and every recentering sequence
$a_R\in\R^4$, $\|a_R\|\to\infty$,
\[
\lim_{R\to\infty}\lim_{N\to\infty}
\int_{\R^8}\Phi(x-a_R,v)\,\pi_{N,1}^{a_R}(dx\,dv)
\to
\int_{\R^8}\Phi(x,v)\,\pi_\infty^0(dx\,dv),
\]
\end{enumerate}

\end{theorem}

\subsection{Imported labeled results from 01\_convergence}

The arguments below use several results from \texttt{01\_convergence} by name.
For later reference, we record the imported statements under the labels used
later in the file.

\begin{theorem}[Imported finite de Finetti theorem on Polish path space]
\label{thm:definetti}
Let \(E\) be a Polish space and let \(P^{(N)}\in\mcP(E^N)\) be exchangeable.
Then for every \(k\le N\) there exists a mixing law
\[
Q_N\in\mcP(\mcP(E))
\]
such that the \(k\)-particle marginal satisfies
\[
d_{\TV}\!\left(
P^{(N)}_k,\;
\int_{\mcP(E)}\eta^{\otimes k}\,Q_N(d\eta)
\right)
\le
\frac{k(k-1)}{2N}.
\]
In \texttt{01\_convergence} this is the Diaconis--Freedman finite de Finetti
theorem applied on standard Borel state spaces.
\end{theorem}

\begin{proposition}[Imported finite-dimensional quadratic moment bounds]
\label{prop:moments}
Under Assumption~\ref{ass:01-core}, the \(N\)-particle kinetic system of
\texttt{01\_convergence} has uniform finite-dimensional moment bounds on every
finite time horizon. In particular, for every \(S>0\) there exists
\(C_S<\infty\), independent of \(N\), such that
\[
\sup_{0\le t\le S}\E\bigl[\|X_1^{(N)}(t)\|^2+\|V_1^{(N)}(t)\|^2\bigr]
\le C_S,
\]
and the corresponding fourth-moment increment bounds needed for
Kolmogorov--Chentsov tightness hold on \(C([0,S];\R^8)\).
\end{proposition}

\begin{lemma}[Imported continuity of the mean-field force on moment classes]
\label{lem:mf-lipschitz}
Let \(F[\mu]\) denote the McKean--Vlasov viscous force field from
\texttt{01\_convergence}. On every class of probability measures on \(\R^8\)
with uniformly bounded second moments, the map
\[
\mu\longmapsto F[\mu]
\]
is continuous in the topology used in the mean-field well-posedness argument;
in particular it is stable under the weak convergence and moment bounds used in
the path-space propagation-of-chaos proof below.
\end{lemma}

\begin{proposition}[Imported McKean--Vlasov well-posedness]
\label{prop:wp}
Under Assumption~\ref{ass:01-core}, the McKean--Vlasov equation of
\texttt{01\_convergence} is well posed on \(\mcP_2(\R^8)\). In particular:
\begin{enumerate}[leftmargin=1.5em,label=\textnormal{(\roman*)}]
\item for every initial law in \(\mcP_2(\R^8)\) there exists a weak solution on
every finite horizon;
\item that weak solution is unique in law.
\end{enumerate}
\end{proposition}

\begin{proposition}[Imported stationary mean-field existence and uniqueness]
\label{prop:mf-unique}
Under Assumption~\ref{ass:01-core}, there exists a unique stationary
McKean--Vlasov law
\[
\pi_\infty\in\mcP(\R^8)
\]
with finite second moments. Writing \(\rho_\infty\) for its density and
\(\mathcal L_\mu\) for the mean-field kinetic generator
\[
\mathcal L_\mu \phi(x,v)
:=
v\cdot\nabla_x\phi(x,v)
+\bigl(-\nabla U(x)-\gamma v+F[\mu](x,v)\bigr)\cdot\nabla_v\phi(x,v)
+\frac{\sigma^2}{2}\Delta_v\phi(x,v),
\]
the stationary law satisfies the weak stationary equation
\begin{equation}
\label{eq:mf-stationary-weak}
\int_{\R^8}\mathcal L_{\pi_\infty}\phi(x,v)\,\pi_\infty(dx\,dv)=0
\qquad
\text{for every }\phi\in C_c^\infty(\R^8).
\end{equation}
\end{proposition}

\subsection{Thermodynamic transport of graph and plaquette observables}

The next results transfer empirical trajectory information to the interaction
graph and its elementary plaquettes.

\begin{definition}[One-particle state and empirical stationary measure]
\label{def:one-particle-state}
For each \(N\), particle index \(i\), and time slice \(t\in\Tset\), define the
one-particle phase-space state
\[
Y_i^{(N)}(t):=\bigl(X_i^{(N)}(t),V_i^{(N)}(t)\bigr)\in\R^8,
\]
and the empirical same-time measure
\[
L_t^{(N)}
:=
\frac1N\sum_{i=1}^N \delta_{Y_i^{(N)}(t)}.
\]
In the notation of Definition~\ref{def:bridge-notation}, \(L_t^{(N)}=\mu_t^N\).
For the path-space statements below, we use the same symbol
\[
Y_i^{(N)}(s)\in\R^8,\qquad s\in[0,S],
\]
for the continuous-time kinetic path of particle \(i\) under the
\(N\)-particle dynamics of \texttt{01\_convergence}, started from the invariant
law \(\pi_N\).
The recorded discrete data are the observed values of this path at the
integer-indexed sampling times.
\end{definition}

\begin{corollary}[Empirical-measure collapse at stationarity]
\label{cor:stationary-empirical-collapse}
Under Assumption~\ref{ass:01-core}, for every fixed recorded time \(t\in\Tset\),
the empirical measure \(L_t^{(N)}\) converges in probability to the unique
stationary mean-field law \(\pi_\infty\):
\[
L_t^{(N)}
\xrightarrow[N\to\infty]{\mathbb P}
\pi_\infty
\qquad\text{in }\mcP(\R^8).
\]
\end{corollary}

\begin{proof}
At \(t=0\), this is exactly the collapse of the de Finetti mixing law proved in
the proof of the stationary mean-field limit theorem imported from
\texttt{01\_convergence}: with the canonical choice
\[
Q_N:=\Law\!\left(\frac1N\sum_{i=1}^N\delta_{Y_i^{(N)}(0)}\right),
\]
the tightness--identification--uniqueness argument shows \(Q_N\Rightarrow
\delta_{\pi_\infty}\).
Equivalently,
\[
\frac1N\sum_{i=1}^N\delta_{Y_i^{(N)}(0)}
\xrightarrow[N\to\infty]{\mathbb P}\pi_\infty.
\]
Because the system is started from its invariant law \(\pi_N\), the law of
\(L_t^{(N)}\) does not depend on \(t\), so the same convergence holds for every
fixed recorded slice \(t\in\Tset\).
\end{proof}

\begin{definition}[Label-invariant pair rule and empirical edge measure]
\label{def:edge-measure}
Fix a recorded time slice \(t\in\Tset\).
Define the same-time companion edge set by
\[
\mcP_t^{(N)}
:=
\{(i,j)\in\{1,\dots,N\}^2:i\neq j\}.
\]
Assume that the same-time companion relation is generated by a label-invariant
pair rule, in the sense that there exists a bounded measurable indicator
\[
\chi_t:\R^8\times\R^8\to\{0,1\}
\]
such that for all \(N\) and all \(i\neq j\),
\[
\one_{\{(i,j)\in\mcP_t^{(N)}\}}
=
\chi_t\!\bigl(Y_i^{(N)}(t),Y_j^{(N)}(t)\bigr).
\]
Define the empirical edge measure by
\[
\Gamma_t^{(N)}
:=
\frac1{N^2}\sum_{i\neq j}
\one_{\{(i,j)\in\mcP_t^{(N)}\}}
\delta_{\bigl(Y_i^{(N)}(t),Y_j^{(N)}(t)\bigr)}.
\]
\end{definition}

\begin{lemma}[Weak convergence against bounded a.e.-continuous test functions]
\label{lem:weak-ae-continuous}
Let \(\mu_N\Rightarrow\mu\) in \(\mcP(E)\) and \(f:E\to\R\) be bounded with
discontinuity set \(\operatorname{Disc}(f)\). If
\(\mu(\operatorname{Disc}(f))=0\), then
\[
\int_E f\,d\mu_N\to \int_E f\,d\mu.
\]
This is the Portmanteau theorem for bounded functions with \(\mu\)-a.e.
continuity, see for example
\cite[Theorem 2.1]{billingsley1999convergence}
(the Portmanteau theorem).
\end{lemma}

\begin{theorem}[Stationary thermodynamic limit of the interaction graph]
\label{thm:edge-measure-limit}
Assume Definition~\ref{def:edge-measure}.
Let \(\Phi:\R^8\times\R^8\to\R\) be bounded and continuous, and assume that
\(\chi_t\Phi\) is continuous at \((\pi_\infty\otimes\pi_\infty)\)-almost every
pair.
Then
\[
\int_{\R^8\times\R^8}\Phi\,d\Gamma_t^{(N)}
\xrightarrow[N\to\infty]{\mathbb P}
\iint_{\R^8\times\R^8}
\chi_t(z,z')\,\Phi(z,z')\,\pi_\infty(dz)\,\pi_\infty(dz').
\]
Equivalently,
\[
\frac1{N^2}\sum_{(i,j)\in\mcP_t^{(N)}}
\Phi\!\bigl(Y_i^{(N)}(t),Y_j^{(N)}(t)\bigr)
\xrightarrow[N\to\infty]{\mathbb P}
\iint_{\R^8\times\R^8}
\chi_t(z,z')\,\Phi(z,z')\,\pi_\infty(dz)\,\pi_\infty(dz').
\]
If the number of recorded slices \(T+1\) is fixed, the same convergence holds
after summing over finitely many \(t\in\Tset\).
\end{theorem}

\begin{proof}
Write
\[
\Psi_t(z,z'):=\chi_t(z,z')\,\Phi(z,z').
\]
By Definition~\ref{def:edge-measure},
\begin{align*}
\int \Phi\,d\Gamma_t^{(N)}
&=
\frac1{N^2}\sum_{i\neq j}\Psi_t\!\bigl(Y_i^{(N)}(t),Y_j^{(N)}(t)\bigr)\\
&=
\frac1{N^2}\sum_{i,j}\Psi_t\!\bigl(Y_i^{(N)}(t),Y_j^{(N)}(t)\bigr)
-\frac1{N^2}\sum_{i=1}^N\Psi_t\!\bigl(Y_i^{(N)}(t),Y_i^{(N)}(t)\bigr)\\
&=
\iint_{\R^8\times\R^8}\Psi_t(z,z')\,L_t^{(N)}(dz)\,L_t^{(N)}(dz')
-R_t^{(N)},
\end{align*}
where
\[
|R_t^{(N)}|
\le
\frac{\|\Psi_t\|_\infty}{N}.
\]
Hence \(R_t^{(N)}\to0\) in probability.

By Corollary~\ref{cor:stationary-empirical-collapse},
\[
L_t^{(N)}\xrightarrow{\mathbb P}\pi_\infty.
\]
Since the limit is deterministic, the product measure
\(L_t^{(N)}\otimes L_t^{(N)}\) converges in probability to
\(\pi_\infty\otimes\pi_\infty\).
Since \((\pi_\infty\otimes\pi_\infty)(\operatorname{Disc}\Psi_t)=0\),
where \(\operatorname{Disc}\Psi_t\) is the discontinuity set, Lemma
\ref{lem:weak-ae-continuous} yields
\[
\iint \Psi_t(z,z')\,L_t^{(N)}(dz)\,L_t^{(N)}(dz')
\xrightarrow[N\to\infty]{\mathbb P}
\iint \Psi_t(z,z')\,\pi_\infty(dz)\,\pi_\infty(dz').
\]
Combining the last two displays proves the first claim.
The second is the same identity written in index notation.
If \(T+1\) is fixed, summing over finitely many times preserves convergence in
probability.
\end{proof}

\begin{definition}[Adjacent-time empirical path and plaquette measures]
\label{def:plaquette-measure}
Fix \(0\le t<T\) and define
\[
\widetilde Y_i^{(N)}(t)
:=
\bigl(Y_i^{(N)}(t),Y_i^{(N)}(t+1)\bigr)\in\R^{12},
\qquad
M_t^{(N)}
:=
\frac1N\sum_{i=1}^N \delta_{\widetilde Y_i^{(N)}(t)}.
\]
Define the persistent-plaquette indicator by
\[
\chi_{\square,t}\bigl((z_0,z_1),(w_0,w_1)\bigr)
:=
\chi_t(z_0,w_0)\,\chi_{t+1}(z_1,w_1),
\]
and the empirical plaquette measure by
\[
\Gamma_{\square,t}^{(N)}
:=
\frac1{N^2}\sum_{i\neq j}
\chi_{\square,t}\!\bigl(\widetilde Y_i^{(N)}(t),\widetilde Y_j^{(N)}(t)\bigr)\,
\delta_{\bigl(\widetilde Y_i^{(N)}(t),\widetilde Y_j^{(N)}(t)\bigr)}.
\]
\end{definition}

\begin{definition}[Stationary mean-field path law]
\label{def:stationary-path-law}
For every finite horizon \(S>0\), let \(\Pi_{[0,S]}^\infty\) be the law on
\(C([0,S];\R^8)\) of the unique McKean--Vlasov solution with initial law
\(\pi_\infty\).
Because \(\pi_\infty\) is stationary, \(\Pi_{[0,S]}^\infty\) is stationary under
time shifts.
For each integer \(0\le t\le\lfloor S\rfloor-1\), define the adjacent-time mean-field law by
\[
\Lambda_t^\infty
:=
(\operatorname{ev}_{t,t+1})_\#\Pi_{[0,S]}^\infty,
\qquad
\operatorname{ev}_{t,t+1}(\gamma):=\bigl(\gamma(t),\gamma(t+1)\bigr).
\]
\end{definition}

\begin{definition}[Finite-horizon shift-stationary path laws]
\label{def:shift-stationary-path-laws}
For finite horizon \(S>0\), set
\[
\mathcal S^{\mathrm{stat}}_{[0,S]}
:=
\bigcap_{t=0}^{\lfloor S\rfloor-1}
\bigl\{
\eta\in\mcP(C([0,S];\R^8)):
(\operatorname{ev}_{t,t+1})_\#\eta
=
(\operatorname{ev}_{0,1})_\#\eta
\bigr\}.
\]
This is the class of one-particle path laws whose unit-time adjacent marginals on
the discrete times \(0,\dots,\lfloor S\rfloor-1\) are stationary.
\end{definition}

\begin{lemma}[Shift-stationary class is closed under weak limits]
\label{lem:shift-stationary-closed}
Let \(E:=C([0,S];\R^8)\). The set
\[
\mathcal S^{\mathrm{stat}}_{[0,S]}\subset\mcP(E)
\]
is Borel (indeed, closed). If \(\eta_n\Rightarrow\eta\) in \(\mcP(E)\) and
\(\eta_n\in\mathcal S^{\mathrm{stat}}_{[0,S]}\) for all \(n\), then
\(\eta\in\mathcal S^{\mathrm{stat}}_{[0,S]}\).
Moreover, if \(Q_n\Rightarrow Q\) in \(\mcP(\mcP(E))\) and
\(Q_n(\mathcal S^{\mathrm{stat}}_{[0,S]})=1\) for all \(n\), then
\(Q(\mathcal S^{\mathrm{stat}}_{[0,S]})=1\).
\end{lemma}

\begin{proof}
For each \(t\in\{0,\dots,\lfloor S\rfloor-1\}\), the map
\[
\eta\mapsto(\operatorname{ev}_{t,t+1})_\#\eta
\]
from \(\mcP(E)\) to \(\mcP(\R^8\times\R^8)\) is continuous under weak
convergence of \(\mcP(E)\). Hence
\[
\eta\mapsto\left(
(\operatorname{ev}_{t,t+1})_\#\eta,\,
(\operatorname{ev}_{0,1})_\#\eta
\right)
\]
is continuous, so each set
\[
\{\eta:(\operatorname{ev}_{t,t+1})_\#\eta=(\operatorname{ev}_{0,1})_\#\eta\}
\]
is closed as a preimage of the diagonal in
\(\mcP(\R^8\times\R^8)\times\mcP(\R^8\times\R^8)\). Finite intersection
gives that \(\mathcal S^{\mathrm{stat}}_{[0,S]}\) is closed. The final
statement for \(Q_n\) is the Portmanteau theorem applied to this closed set.
\end{proof}

\begin{theorem}[Stationary path-space propagation of chaos]
\label{thm:stationary-path-poc}
Under Assumption~\ref{ass:01-core}, for every finite horizon \(S>0\), let
\(\widehat P_{[0,S]}^{(N)}\) be the law on
\(C([0,S];(\R^8)^N)\) of the \(N\)-particle kinetic system started from its
invariant law \(\pi_N\), and define the empirical path measure
\[
\mathfrak M_{[0,S]}^{(N)}
:=
\frac1N\sum_{i=1}^N \delta_{Y_i^{(N)}(\cdot)}
\in\mcP\bigl(C([0,S];\R^8)\bigr).
\]
Then
\[
\mathfrak M_{[0,S]}^{(N)}
\xrightarrow[N\to\infty]{\mathbb P}
\Pi_{[0,S]}^\infty
\qquad
\text{in }\mcP\bigl(C([0,S];\R^8)\bigr).
\]
Equivalently, for every bounded continuous path functional
\(\Psi:C([0,S];\R^8)\to\R\),
\[
\frac1N\sum_{i=1}^N \Psi\!\bigl(Y_i^{(N)}(\cdot)\bigr)
\xrightarrow[N\to\infty]{\mathbb P}
\int_{C([0,S];\R^8)}\Psi(\gamma)\,\Pi_{[0,S]}^\infty(d\gamma).
\]
In particular, for each integer \(0\le t\le\lfloor S\rfloor-1\),
\[
M_t^{(N)}
=(\operatorname{ev}_{t,t+1})_\# \mathfrak M_{[0,S]}^{(N)}
\xrightarrow[N\to\infty]{\mathbb P}
\Lambda_t^\infty.
\]
\end{theorem}

\begin{proof}
Write \(E:=C([0,S];\R^8)\).
Because the invariant law \(\pi_N\) is exchangeable and the \(N\)-particle
dynamics are permutation-equivariant, the stationary path law
\(\widehat P_{[0,S]}^{(N)}\) on \(E^N\) is exchangeable.
Let
\[
\mathfrak M_{[0,S]}^{(N)}(\omega)
:=
\frac1N\sum_{i=1}^N\delta_{Y_i^{(N)}(\cdot)(\omega)}
\in\mcP(E)
\]
be the empirical measure of particle trajectories under
\(\widehat P_{[0,S]}^{(N)}\). Since \(E\) is Polish, hence standard Borel,
Theorem~\ref{thm:definetti} from \texttt{01\_convergence}, i.e.\ the
Diaconis--Freedman finite de Finetti theorem of
\cite{diaconis1980finite},
applies directly on the path space \(E=C([0,S];\R^8)\), yielding for every
\(k\le N\) a mixing law
\(\widehat Q_N\in\mcP(\mcP(E))\) such that
\[
\widehat Q_N:=\Law\!\bigl(\mathfrak M_{[0,S]}^{(N)}\bigr)
\in\mcP(\mcP(E)),
\]
and
\[
d_{\TV}\!\left(
\widehat P^{(N)}_{[0,S],k},\;
\int_{\mcP(E)}\eta^{\otimes k}\,\widehat Q_N(d\eta)
\right)
\le
\frac{k(k-1)}{2N}.
\]
Hence the fixed-\(k\) path marginals of \(\widehat P_{[0,S]}^{(N)}\) are within
\(O(1/N)\) in total variation of mixtures of product path laws.
In particular, for \(k=1\), \(\widehat P^{(N)}_{[0,S],1}\in\mathcal S^{\mathrm{stat}}_{[0,S]}\),
and Lemma~\ref{lem:shift-stationary-closed} applies to weak limits of such
stationary one-particle path laws. This records that adjacent-time
stationarity is stable under the one-particle weak-limit passage. For the
second-level laws \(\widehat Q_N\), however, we do not use any support claim on
\(\mathcal S^{\mathrm{stat}}_{[0,S]}\); the identification of subsequential
limits below proceeds instead through the nonlinear martingale problem and the
uniqueness of the McKean--Vlasov path law for the initial distribution
\(\pi_\infty\).

The one-particle path marginals are tight on \(E\).
Indeed, stationarity gives the same uniform moment bounds at every time as in the
bridge package from \texttt{01\_convergence} (in particular
Assumption~\ref{ass:bridge-backbone}(ii) and (iii)), and
Proposition~\ref{prop:moments} of
\texttt{01\_convergence} (the finite-dimensional quadratic moment bounds) provides
the required uniform-in-\(N\) bounds.
In particular, combining the linear-growth drift, additive velocity noise, and
Burkholder--Davis--Gundy estimates for the velocity increments gives
\[
\E\|Y_1^{(N)}(t)-Y_1^{(N)}(s)\|^4
\le C|t-s|^2,
\]
uniformly in \(N\), which implies tightness of
\(\Law(Y_1^{(N)}(\cdot))\) in \(E\) by the Kolmogorov--Chentsov criterion
(see \cite[Chapter~2]{billingsley1999convergence}). Hence tightness of the second-level laws
\(\{\widehat Q_N\}_{N\ge2}\). Since \(E\) is Polish, \(\mcP(E)\) is Polish under
the weak topology, and therefore \(\{\widehat Q_N\}_{N\ge2}\) is tight in
\(\mcP(E)\).
Let \(\widehat Q_\infty\) be any subsequential limit.

We identify \(\widehat Q_\infty\)-a.e. limit paths by the nonlinear martingale
problem.
Fix \(\phi\in C_c^\infty(\R^8)\), times \(0\le u\le v\le S\), and a bounded
continuous functional \(G\) depending only on the path segment on \([0,u]\).
For \(\mu\in\mcP(\R^8)\), write
\[
\mathcal L_\mu \phi(x,v)
:=
v\cdot\nabla_x\phi(x,v)
+\bigl(-\nabla U(x)-\gamma v+F[\mu](x,v)\bigr)\cdot\nabla_v\phi(x,v)
+\frac{\sigma^2}{2}\Delta_v\phi(x,v),
\]
where \(F[\mu]\) is the mean-field viscous field from
\texttt{01\_convergence}.
For the finite-\(N\) system,
\begin{multline*}
\phi(Y_i^{(N)}(v))-\phi(Y_i^{(N)}(u))
-\int_u^v
\Bigl[
\mathcal L_{\mathfrak M_{[0,S],s}^{(N)}}\phi\bigl(Y_i^{(N)}(s)\bigr)\\
+\bigl(
F_i^{(N)}(s)-F[\mathfrak M_{[0,S],s}^{(N)}]\bigl(Y_i^{(N)}(s)\bigr)
\bigr)
\cdot\nabla_v\phi\bigl(Y_i^{(N)}(s)\bigr)
\Bigr]\,ds
\end{multline*}
is a martingale, where
\[
\mathfrak M_{[0,S],s}^{(N)}
:=
(\operatorname{ev}_s)_\#\mathfrak M_{[0,S]}^{(N)}
=
\frac1N\sum_{j=1}^N \delta_{Y_j^{(N)}(s)}.
\]
Averaging over \(i\) and multiplying by \(G(Y_i^{(N)}|_{[0,u]})\) gives zero
expectation.
The force-defect estimate from \texttt{01\_convergence}, integrated on
\([u,v]\), shows that the defect term contributes only \(O(N^{-1/2})\) in
expectation.
Passing to the limit and using the continuity of \(\mu\mapsto F[\mu]\) on
uniformly moment-bounded classes (Lemma~\ref{lem:mf-lipschitz} in
\texttt{01\_convergence}, with the moment class provided by
Assumption~\ref{ass:bridge-backbone}(ii)), together with continuity of
\(\eta\mapsto\eta_s=(\operatorname{ev}_s)_\#\eta\), \(\widehat Q_\infty\)-almost
every path law \(\eta\) satisfies
\[
\int_E
G(\gamma|_{[0,u]})
\Bigl[
\phi(\gamma(v))-\phi(\gamma(u))
-\int_u^v \mathcal L_{\eta_s}\phi(\gamma(s))\,ds
\Bigr]
\eta(d\gamma)
=
0,
\]
where \(\eta_s:=(\operatorname{ev}_s)_\#\eta\).
Thus \(\widehat Q_\infty\)-almost every \(\eta\) solves the McKean--Vlasov
martingale problem on \([0,S]\).

Because each \(\widehat P_{[0,S]}^{(N)}\) is stationary in time (as the path
law induced by the invariant measure \(\pi_N\)), stationarity of time marginals
implies, for every \(N\),
\[
(\operatorname{ev}_0)_\#\widehat Q_N
\;=\;
\Law\!\left((\operatorname{ev}_0)_\#\mathfrak M_{[0,S]}^{(N)}\right)
\;=\;
\Law(L_0^{(N)}).
\]
By Corollary~\ref{cor:stationary-empirical-collapse} with \(t=0\),
\(L_0^{(N)}\xrightarrow{\mathbb P}\pi_\infty\), hence
\((\operatorname{ev}_0)_\#\widehat Q_N\Rightarrow\delta_{\pi_\infty}\)
in \(\mcP(\R^8)\). Since \(\mathrm{ev}_0:E\to\R^8\) is continuous, the
pushforward map \(\eta\mapsto(\operatorname{ev}_0)_\#\eta\) is continuous on
\(\mcP(E)\) in the weak topology (continuous mapping theorem). This gives
\[
(\operatorname{ev}_0)_\#\widehat Q_\infty=\delta_{\pi_\infty},
\]
so \(\widehat Q_\infty\)-almost every \(\eta\) satisfies \(\eta_0=\pi_\infty\).
Assumption~\ref{ass:bridge-backbone}(iii) and weak convergence
\(\pi_{N,1}\Rightarrow\pi_\infty\) imply \(\pi_\infty\in\mcP_2(\R^8)\), so
we may apply this well-posedness statement with initial law \(\pi_\infty\):
By Proposition~\ref{prop:wp}(ii) of \texttt{01\_convergence}, the McKean--Vlasov
equation from that framework is well posed for every initial law in
\(\mcP_2(\R^8)\), hence there is a unique path law with initial law
\(\pi_\infty\); by Definition~\ref{def:stationary-path-law}, that unique law is
\(\Pi_{[0,S]}^\infty\).
Hence
\[
\widehat Q_\infty=\delta_{\Pi_{[0,S]}^\infty}.
\]
Every subsequence has the same limit, so
\[
\Law\!\bigl(\mathfrak M_{[0,S]}^{(N)}\bigr)
\Rightarrow
\delta_{\Pi_{[0,S]}^\infty},
\]
which is equivalent to the stated convergence in probability.
The final claim follows from the continuous mapping theorem for the evaluation
map \(\operatorname{ev}_{t,t+1}\). As \(\operatorname{ev}_{t,t+1}:E\to
\R^8\times\R^8\) is continuous, this gives the claimed pushforward limit.
\end{proof}

\begin{theorem}[Stationary thermodynamic limit of plaquette observables]
\label{thm:plaquette-measure-limit}
Let \(\Phi:(\R^{12})^2\to\R\) be bounded and continuous, and assume that
\(\chi_{\square,t}\Phi\) is continuous at
\((\Lambda_t^\infty\otimes\Lambda_t^\infty)\)-almost every pair.
Then
\[
\int_{(\R^{12})^2}\Phi\,d\Gamma_{\square,t}^{(N)}
\xrightarrow[N\to\infty]{\mathbb P}
\iint_{(\R^{12})^2}
\chi_{\square,t}(\zeta,\zeta')\,\Phi(\zeta,\zeta')\,
\Lambda_t^\infty(d\zeta)\,\Lambda_t^\infty(d\zeta').
\]
\end{theorem}

\begin{proof}
Set
\[
\Psi_{\square,t}(\zeta,\zeta')
:=
\chi_{\square,t}(\zeta,\zeta')\,\Phi(\zeta,\zeta').
\]
By Definition~\ref{def:plaquette-measure},
\begin{align*}
\int \Phi\,d\Gamma_{\square,t}^{(N)}
&=
\frac1{N^2}\sum_{i\neq j}
\Psi_{\square,t}\!\bigl(\widetilde Y_i^{(N)}(t),\widetilde Y_j^{(N)}(t)\bigr)\\
&=
\iint \Psi_{\square,t}(\zeta,\zeta')\,M_t^{(N)}(d\zeta)\,M_t^{(N)}(d\zeta')
\;-\;R_{\square,t}^{(N)},
\end{align*}
where
\[
|R_{\square,t}^{(N)}|
\le
\frac{\|\Psi_{\square,t}\|_\infty}{N}.
\]
Hence \(R_{\square,t}^{(N)}\to0\) in probability.
Because
\[
M_t^{(N)}
\xrightarrow[N\to\infty]{\mathbb P}
\Lambda_t^\infty
\]
by Theorem~\ref{thm:stationary-path-poc}, the product measure
\(M_t^{(N)}\otimes M_t^{(N)}\) converges in probability to
\(\Lambda_t^\infty\otimes\Lambda_t^\infty\).
Since \((\Lambda_t^\infty\otimes\Lambda_t^\infty)(\operatorname{Disc}
\Psi_{\square,t})=0\), where \(\operatorname{Disc}\Psi_{\square,t}\) is the
discontinuity set, Lemma~\ref{lem:weak-ae-continuous} yields
\[
\iint \Psi_{\square,t}(\zeta,\zeta')\,M_t^{(N)}(d\zeta)\,M_t^{(N)}(d\zeta')
\xrightarrow[N\to\infty]{\mathbb P}
\iint \Psi_{\square,t}(\zeta,\zeta')\,
\Lambda_t^\infty(d\zeta)\,\Lambda_t^\infty(d\zeta').
\]
Combining the last two displays proves the theorem.
\end{proof}

\section{Continuum comparison field}

\begin{definition}[Continuum color field and induced gauge curvature]
\label{def:continuum-color-field}
Let
\[
c:\R^4\to\C^3
\]
be a differentiable color field. Its induced continuum connection is the
\(\mathfrak h_3\)-valued one-form
\[
A_\mu[c](x)
:=
\frac12\sum_{a=1}^8
\Re\!\bigl(c(x)^\dagger\lambda_a\,\partial_\mu c(x)\bigr)\lambda_a,
\qquad
0\le \mu\le 3,
\]
and its induced continuum curvature is
\[
F_{\mu\nu}[c]
:=
\partial_\mu A_\nu[c]-\partial_\nu A_\mu[c]
+ig_3[A_\mu[c],A_\nu[c]].
\]
\end{definition}

\begin{remark}[Convention used in this paper]
\label{rem:hermitian-convention}
The formulas in this paper are written in the Hermitian generator convention:
\(A_\mu[c]\) and \(F_{\mu\nu}[c]\) take values in the Hermitian traceless space
\(\mathfrakh\), and the corresponding \(SU(3)\) transports are written as
exponentials \(\exp(iH)\) with \(H\in\mathfrakh\).
Proposition~\ref{prop:su3-convention-bridge} below translates these objects to
the standard anti-Hermitian \(\mathfrak{su}(3)\) convention used in the usual
statement of pure \(SU(3)\) Yang--Mills theory.
\end{remark}

\section{Realized Fractal regulator}

The bridge package imported from \texttt{01\_convergence} supplies the
stochastic substrate. In the revised formulation, that substrate is not sampled
onto an externally imposed lattice. Instead it is pushed directly into the
Fractal regulator extracted from the run history itself.

\begin{definition}[Recorded swarm history and finite path law]
\label{def:history-path-law}
Fix \(N\ge N_0\) and a finite recorded horizon \(T\in\N\). Let
\[
\Omega_{N,T}
\]
denote the measurable space of recorded swarm histories
\[
\omega=
\bigl(
Z^{(N)}(0),\dots,Z^{(N)}(T),
\mcC^{(N)},\mcB^{(N)},\mcR^{(N)}
\bigr),
\]
where \(Z^{(N)}(t)\in(\R^8)^N\) is the recorded particle state at time \(t\),
\(\mcC^{(N)}\) collects the realized companion selections, \(\mcB^{(N)}\)
collects the realized cloning/birth-death decisions, and \(\mcR^{(N)}\)
collects the realized attribution data and edge-rotation payloads produced by
the run. The corresponding finite swarm path law is
\[
\mathbb P_{N,T}:=\Law(\omega)\in\mcP(\Omega_{N,T}).
\]
\end{definition}

\begin{definition}[Realized Fractal regulator]
\label{def:realized-fractal-regulator}
For each history \(\omega\in\Omega_{N,T}\), let
\[
\mathfrak F^{(N,T)}(\omega)
:=
\bigl(
\mcV^{(N,T)}(\omega),
\mcE_{\mathrm{CST}}^{(N,T)}(\omega),
\mcE_{\mathrm{IG}}^{(N,T)}(\omega),
\mcE_{\mathrm{IA}}^{(N,T)}(\omega),
\mcP_{\mathrm{FG}}^{(N,T)}(\omega),
\Xi^{(N,T)}(\omega)
\bigr)
\]
be the realized Fractal regulator extracted from the run history, with the
following components.
\begin{enumerate}[leftmargin=1.5em,label=\textnormal{(\roman*)}]
\item \textbf{Nodes.}
\(\mcV^{(N,T)}(\omega)\) consists of the recorded pre/clone/final events
\[
n_{i,s}^{\mathrm{pre}},\qquad
n_{i,s}^{\mathrm{clone}},\qquad
n_{i,s}^{\mathrm{fin}},
\]
for walker \(i\in\{1,\dots,N\}\) and recorded time \(s\in\{0,\dots,T\}\),
restricted to the stages present in the run history.
\item \textbf{CST edges.}
\(\mcE_{\mathrm{CST}}^{(N,T)}(\omega)\) contains the temporal genealogy edges
connecting successive stages of the same walker and the inherited
parent-to-child temporal edges created by cloning events.
\item \textbf{IG edges.}
\(\mcE_{\mathrm{IG}}^{(N,T)}(\omega)\) contains the realized companion edges at
common recorded time, joining each walker to its realized companion according
to \(\mcC^{(N)}\).
\item \textbf{IA edges.}
\(\mcE_{\mathrm{IA}}^{(N,T)}(\omega)\) contains the influence-attribution edges
that close the interaction triangles and carry the non-abelian credit
assignment data extracted from \(\mcR^{(N)}\).
\item \textbf{Interaction plaquettes.}
\(\mcP_{\mathrm{FG}}^{(N,T)}(\omega)\) is the family of oriented interaction
\(4\)-cycles obtained by gluing a pair of opposite interaction triangles along
their common IG edge. Equivalently, the plaquette boundary uses two temporal
segments and two IA back-edges, while the shared IG edge cancels by opposite
orientation.
\item \textbf{Embedding.}
\(\Xi^{(N,T)}(\omega)\) maps each node to its realized spacetime event
\[
\Xi^{(N,T)}(\omega,n)=(\tau_n(\omega),x_n(\omega))\in\R\times\R^4,
\]
where \(\tau_n\) is the recorded time coordinate and \(x_n\) the realized
position.
\end{enumerate}
\end{definition}

\begin{definition}[Canonical plaquette geometry from the realized embedding]
\label{def:fractal-plaquette-geometry}
For each plaquette \(P\in\mcP_{\mathrm{FG}}^{(N,T)}(\omega)\), write its
oriented boundary in the order supplied by the realized Fractal regulator as
\[
\partial P=(v_0\to v_1\to v_2\to v_3\to v_4=v_0).
\]
Using the realized embedding
\[
\Xi^{(N,T)}(\omega,v_r)=\bigl(\tau_{v_r}(\omega),x_{v_r}(\omega)\bigr)\in\R\times\R^4,
\]
define the plaquette base point by the canonical barycenter
\[
x_P(\omega):=\frac14\sum_{r=0}^3 x_{v_r}(\omega)\in\R^4,
\]
and define the ordered edge displacements by
\[
\tau_{P,r}(\omega):=x_{v_{r+1}}(\omega)-x_{v_r}(\omega)\in\R^4,
\qquad 0\le r\le 3.
\]
The effective oriented area bivector is then defined by the polygonal boundary
formula
\[
\Sigma_P^{\alpha\beta}(\omega)
:=
\frac12\sum_{r}
\bigl(
\tau_{P,r}^{\alpha}(\omega)\tau_{P,r+1}^{\beta}(\omega)
-\tau_{P,r}^{\beta}(\omega)\tau_{P,r+1}^{\alpha}(\omega)
\bigr).
\]
The effective plaquette area is
\[
A_P(\omega)
:=
\Bigl(
\frac12\sum_{\alpha,\beta=0}^3
\Sigma_P^{\alpha\beta}(\omega)\Sigma_P^{\alpha\beta}(\omega)
\Bigr)^{1/2}.
\]
All quantities in this definition depend only on the realized oriented boundary
vertices of \(P\) and the ambient coordinates \(x_v(\omega)\); no auxiliary
chart choice enters.
\end{definition}

\begin{proposition}[Measurable finite regulator extracted from the swarm]
\label{prop:fractal-regulator-measurable}
For each \(N,T\), the map
\[
\mathfrak F^{(N,T)}:\Omega_{N,T}\to
\{\text{finite oriented \(2\)-complexes with canonical embedded plaquette data}\}
\]
defined by Definition~\ref{def:realized-fractal-regulator} and
Definition~\ref{def:fractal-plaquette-geometry} is measurable. Moreover, for
every \(\omega\in\Omega_{N,T}\), the sets
\[
\mcV^{(N,T)}(\omega),\quad
\mcE_{\mathrm{CST}}^{(N,T)}(\omega),\quad
\mcE_{\mathrm{IG}}^{(N,T)}(\omega),\quad
\mcE_{\mathrm{IA}}^{(N,T)}(\omega),\quad
\mcP_{\mathrm{FG}}^{(N,T)}(\omega)
\]
are finite.
\end{proposition}

\begin{proof}
The recorded history contains finitely many walkers, finitely many observation
times, finitely many companion choices, finitely many cloning events, and
finitely many attribution payloads. Every node, edge, plaquette, barycenter,
and canonical plaquette quantity used above is therefore obtained by a finite
sequence of measurable selections and algebraic operations from the recorded
history.
This proves measurability and finiteness.
\end{proof}

\section{Stationary mean-field comparison profile}

\begin{definition}[Stationary mean-field color profile]
\label{def:stationary-mf-color}
Let \(\pi_\infty\) be the unique stationary mean-field law from
Proposition~\ref{prop:mf-unique}, and write \(\rho_\infty(x,v)\) for its
density on \(\R^8=\R_x^4\times\R_v^4\). Define the stationary spatial marginal
\[
\rho_{\infty,x}(x):=\int_{\R^4}\rho_\infty(x,v)\,\dd v,
\]
the stationary first velocity moment
\[
j_\infty(x):=\int_{\R^4}v\,\rho_\infty(x,v)\,\dd v\in\R^4,
\]
the stationary kernel mass
\[
m_\infty(x)
:=
\iint_{\R^8} K(x,y)\,\rho_\infty(y,q)\,\dd y\,\dd q,
\]
\[
J_\infty(x)
:=
\iint_{\R^8} K(x,y)\,q\,\rho_\infty(y,q)\,\dd y\,\dd q\in\R^4,
\]
the stationary alignment field
\[
A_\infty(x)
:=
\frac{J_\infty(x)}{m_\infty(x)}\in\R^4,
\]
\[
a_\infty(x):=\Pi_{\mathrm{col}}A_\infty(x)\in\R^3,
\]
and the stationary raw color profile
\[
\tilde{\mathbf c}_\infty(x):=a_\infty(x)\in\C^3.
\]
Let
\[
\mathbf c_\infty(x)
:=
\frac{\tilde{\mathbf c}_\infty(x)}
{\sqrt{\tilde{\mathbf c}_\infty(x)^\dagger\tilde{\mathbf c}_\infty(x)+\epsilon_c^2}}.
\]
\end{definition}

\begin{proposition}[Vanishing stationary mean velocity]
\label{prop:stationary-velocity-centering}
The stationary mean-field law satisfies
\[
\int_{\R^8} v\,\pi_\infty(dx\,dv)=0.
\]
Equivalently,
\[
\int_{\R^4} j_\infty(x)\,\dd x=0.
\]
\end{proposition}

\begin{proposition}[Smoothness of the stationary mean-field color profile]
\label{prop:stationary-mf-color-smooth}
The stationary density \(\rho_\infty\) is \(C^\infty\) and strictly positive on
\(\R^8\). Consequently
\[
A_\infty\in C^\infty(\R^4;\R^4),
\qquad
\mathbf c_\infty\in C^\infty(\R^4;\C^3).
\]
\end{proposition}

\begin{proof}[Proof of Proposition~\ref{prop:stationary-velocity-centering}]
Fix \(j\in\{0,1,2,3\}\). Let \(\eta\in C_c^\infty(\R^8)\) satisfy
\(\eta\equiv1\) on the unit ball and \(\eta\equiv0\) outside the ball of radius
\(2\), and set
\[
\phi_{R,j}(x,v):=x_j\,\eta\!\left(\frac{|x|^2+|v|^2}{R^2}\right).
\]
Then \(\phi_{R,j}\in C_c^\infty(\R^8)\), so the stationary weak equation
\eqref{eq:mf-stationary-weak} gives
\[
\int_{\R^8}\mathcal L_{\pi_\infty}\phi_{R,j}(x,v)\,\pi_\infty(dx\,dv)=0.
\]
Because \(\partial_{x_k}\phi_{R,j}=\delta_{jk}+O((|x|^2+|v|^2)^{1/2}/R)\) and
\(\nabla_v\phi_{R,j}=O((|x|^2+|v|^2)^{1/2}/R)\), one has
\[
\mathcal L_{\pi_\infty}\phi_{R,j}(x,v)=v_j+r_{R,j}(x,v),
\]
where \(r_{R,j}\) is supported on
\(\{R\le (|x|^2+|v|^2)^{1/2}\le 2R\}\) and satisfies
\[
|r_{R,j}(x,v)|\le \frac{C}{R}\bigl(1+|x|^2+|v|^2\bigr)
\]
with \(C\) independent of \(R\). By Theorem~\ref{thm:bridge-from-01}\textnormal{(iii)},
\(\pi_\infty\) has finite second moments, hence
\[
\int_{\R^8}|r_{R,j}|\,d\pi_\infty\to0.
\]
Passing to the limit \(R\to\infty\) therefore gives
\[
\int_{\R^8} v_j\,\pi_\infty(dx\,dv)=0.
\]
Since this holds for each \(j\), the vector identity follows. The identity for
\(j_\infty\) is Fubini's theorem.
\end{proof}

\begin{proof}[Proof of Proposition~\ref{prop:stationary-mf-color-smooth}]
Proposition~\ref{prop:mf-unique} provides a stationary weak solution of the
kinetic Fokker--Planck equation with smooth coefficients. The diffusion vector
fields are the velocity derivatives \(\partial_{v_0},\dots,\partial_{v_3}\),
and their commutators with the transport field generate the spatial derivatives
\(\partial_{x_0},\dots,\partial_{x_3}\). Hörmander hypoellipticity
therefore applies, so the stationary weak solution admits a smooth density
\(\rho_\infty\). The associated hypoelliptic heat kernel is strictly positive,
and invariance forces the stationary density to be strictly positive as well.

Since \(K\) is smooth and Proposition~\ref{prop:mf-unique}
provides the required stationary moment bounds, differentiation under the
integral sign is legitimate in the displayed formulas for \(m_\infty\),
\(J_\infty\), and \(A_\infty\). The strict positivity of the kernel floor gives
\(m_\infty(x)\ge\kappa_0>0\). Since \(\Pi_{\mathrm{col}}\) is linear,
\(\tilde{\mathbf c}_\infty=\Pi_{\mathrm{col}}A_\infty\) is smooth as well.
The final normalization with
\(\epsilon_c>0\) preserves smoothness.
\end{proof}

\begin{definition}[Stationary equilibrium color field on \(\R^4\)]
\label{def:stationary-spacetime-color}
The stationary equilibrium color field is the profile of
Definition~\ref{def:stationary-mf-color} itself:
\[
\bar{\tilde{\mathbf c}}_\infty(x):=\tilde{\mathbf c}_\infty(x),
\qquad
\bar{\mathbf c}_\infty(x):=\mathbf c_\infty(x),
\qquad
x\in\R^4.
\]
\end{definition}

\begin{proposition}[Regularity of the induced continuum connection and curvature]
\label{prop:continuum-connection-regularity}
Let
\[
A_\mu^\infty:=A_\mu[\bar{\mathbf c}_\infty],
\qquad
F_{\mu\nu}^\infty:=F_{\mu\nu}[\bar{\mathbf c}_\infty].
\]
Then
\[
A^\infty\in C^\infty(\R^4;\mathfrakh),
\qquad
F^\infty\in C^\infty(\R^4;\mathfrakh).
\]
\end{proposition}

\begin{proof}
By Proposition~\ref{prop:stationary-mf-color-smooth}, the field
\(\mathbf c_\infty\) is \(C^\infty\) on \(\R^4\), and
\(\bar{\mathbf c}_\infty=\mathbf c_\infty\) by
Definition~\ref{def:stationary-spacetime-color}. The formula
in Definition~\ref{def:continuum-color-field} expresses each \(A_\mu^\infty\)
as a finite linear combination of products of \(\bar{\mathbf c}_\infty\) and
its first derivatives, so \(A^\infty\in C^\infty(\R^4;\mathfrakh)\). The same
definition expresses each \(F_{\mu\nu}^\infty\) as a finite linear combination
of first derivatives of \(A^\infty\) and commutators of its components, hence
\(F^\infty\in C^\infty(\R^4;\mathfrakh)\).
\end{proof}

\begin{proposition}[Infrared and ultraviolet control of the equilibrium gauge profile]
\label{prop:continuum-energy-tail}
Let
\[
G(x):=\exp\!\left(-\frac{|x|^2}{2}\right),
\qquad
A_\mu^\infty:=A_\mu[\bar{\mathbf c}_\infty],
\qquad
F_{\mu\nu}^\infty:=F_{\mu\nu}[\bar{\mathbf c}_\infty].
\]
Then:
\begin{enumerate}[leftmargin=1.5em,label=\textnormal{(\roman*)}]
\item \(A_\infty,\nabla A_\infty,\nabla^2A_\infty\in L^1(\R^4)\cap L^\infty(\R^4)\cap C_0(\R^4)\);
\item \(\mathbf c_\infty,\nabla\mathbf c_\infty,\nabla^2\mathbf c_\infty\in
L^\infty(\R^4)\), and \(\nabla\mathbf c_\infty,\nabla^2\mathbf c_\infty\in
L^1(\R^4)\cap C_0(\R^4)\);
\item \(A^\infty,\nabla A^\infty,F^\infty\in L^1(\R^4)\cap L^\infty(\R^4)\cap C_0(\R^4)\);
\item the continuum Yang--Mills density
\[
e_\infty(x):=
\sum_{0\le\mu<\nu\le3}\tr\!\bigl((F_{\mu\nu}^\infty(x))^2\bigr)
\]
belongs to \(W^{1,1}(\R^4)\cap C_0(\R^4)\), and
\[
\lim_{R\to\infty}\int_{|x|>R}e_\infty(x)\,\dd x=0.
\]
\end{enumerate}
\end{proposition}

\begin{proof}
By Theorem~\ref{thm:bridge-from-01}\textnormal{(iii)}, \(\pi_\infty\) has finite
second moments. Hence \(\rho_{\infty,x}\in L^1(\R^4)\), \(j_\infty\in L^1(\R^4;\R^4)\),
and
\[
\int_{\R^4}|y|^2\rho_{\infty,x}(y)\,\dd y
\;+\;
\int_{\R^4}|j_\infty(y)|\,\dd y
<\infty.
\]
By Proposition~\ref{prop:stationary-velocity-centering},
\(\int_{\R^4}j_\infty(y)\,\dd y=0\). Therefore
\[
m_\infty(x)=\kappa_0+(G*\rho_{\infty,x})(x),
\qquad
J_\infty(x)=(G*j_\infty)(x).
\]
For every multi-index \(\beta\) with \(1\le|\beta|\le2\),
\[
\partial^\beta m_\infty=(\partial^\beta G)*\rho_{\infty,x},
\qquad
\partial^\beta J_\infty=(\partial^\beta G)*j_\infty.
\]
Since each \(\partial^\beta G\) is a Schwartz function, Young's inequality gives
\[
\partial^\beta m_\infty,\partial^\beta J_\infty\in
L^1(\R^4)\cap L^\infty(\R^4),
\]
and the Riemann--Lebesgue lemma gives
\[
\partial^\beta m_\infty,\partial^\beta J_\infty\in C_0(\R^4).
\]
Moreover \(m_\infty(x)\ge\kappa_0>0\) everywhere. Quotient differentiation
therefore yields
\[
A_\infty=\frac{J_\infty}{m_\infty},
\qquad
\nabla A_\infty,\nabla^2A_\infty\in
L^1(\R^4)\cap L^\infty(\R^4)\cap C_0(\R^4),
\]
and since \(J_\infty\in C_0(\R^4)\), also \(A_\infty\in C_0(\R^4)\cap L^\infty(\R^4)\).

Because \(\tilde{\mathbf c}_\infty=\Pi_{\mathrm{col}}A_\infty\), the same bounds
hold for \(\tilde{\mathbf c}_\infty\). The map
\[
u\mapsto \frac{u}{\sqrt{u^\dagger u+\epsilon_c^2}}
\]
is a smooth map \(\C^3\to\C^3\) with bounded first and second derivatives, so
the chain rule gives the asserted bounds for \(\mathbf c_\infty\),
\(\nabla\mathbf c_\infty\), and \(\nabla^2\mathbf c_\infty\).

Definition~\ref{def:continuum-color-field} expresses each \(A_\mu^\infty\) as a
finite linear combination of products of \(\bar{\mathbf c}_\infty\) and
\(\partial_\mu\bar{\mathbf c}_\infty\), and each derivative \(\partial_\nu
A_\mu^\infty\) as a finite linear combination of products of
\(\bar{\mathbf c}_\infty\), \(\nabla\bar{\mathbf c}_\infty\), and
\(\nabla^2\bar{\mathbf c}_\infty\). Since \(\bar{\mathbf c}_\infty=\mathbf c_\infty\),
the previous bounds imply
\[
A^\infty,\nabla A^\infty\in L^1(\R^4)\cap L^\infty(\R^4)\cap C_0(\R^4).
\]
The curvature formula
\[
F_{\mu\nu}^\infty
=
\partial_\mu A_\nu^\infty-\partial_\nu A_\mu^\infty+ig_3[A_\mu^\infty,A_\nu^\infty]
\]
then shows \(F^\infty\in L^1(\R^4)\cap L^\infty(\R^4)\cap C_0(\R^4)\).

Repeating the same convolution argument with higher derivatives of \(G\) gives
that the first three spatial derivatives of \(A_\infty\) and \(\mathbf c_\infty\)
belong to \(L^1(\R^4)\cap L^\infty(\R^4)\cap C_0(\R^4)\). Consequently
\(\nabla F^\infty\in L^1(\R^4)\). Finally, each summand
\(\tr((F_{\mu\nu}^\infty)^2)\) is continuous and bounded by
\(C\|F_{\mu\nu}^\infty\|^2\). Because \(F_{\mu\nu}^\infty\in L^1\cap L^\infty\),
its square belongs to \(L^1\), hence \(e_\infty\in L^1(\R^4)\cap C_0(\R^4)\).
Moreover
\[
\nabla e_\infty
=
2\sum_{0\le\mu<\nu\le3}\tr\!\bigl(F_{\mu\nu}^\infty\,\nabla F_{\mu\nu}^\infty\bigr)
\in L^1(\R^4),
\]
so \(e_\infty\in W^{1,1}(\R^4)\).
The vanishing of the tails is the absolute continuity of the integral.
\end{proof}

\section{Realized gauge sector and induced gauge law}

\begin{definition}[Gell-Mann basis and Hermitian generator space]
\label{def:gellmann}
The Gell-Mann matrices are
\[
\lambda_1=\begin{pmatrix}0&1&0\\1&0&0\\0&0&0\end{pmatrix},\quad
\lambda_2=\begin{pmatrix}0&-i&0\\ i&0&0\\0&0&0\end{pmatrix},\quad
\lambda_3=\begin{pmatrix}1&0&0\\0&-1&0\\0&0&0\end{pmatrix},
\]
\[
\lambda_4=\begin{pmatrix}0&0&1\\0&0&0\\1&0&0\end{pmatrix},\quad
\lambda_5=\begin{pmatrix}0&0&-i\\0&0&0\\ i&0&0\end{pmatrix},\quad
\lambda_6=\begin{pmatrix}0&0&0\\0&0&1\\0&1&0\end{pmatrix},
\]
\[
\lambda_7=\begin{pmatrix}0&0&0\\0&0&-i\\0&i&0\end{pmatrix},\quad
\lambda_8=\frac{1}{\sqrt3}
\begin{pmatrix}1&0&0\\0&1&0\\0&0&-2\end{pmatrix}.
\]
Each $\lambda_a$ is Hermitian and traceless, and
\[
\tr(\lambda_a\lambda_b)=2\delta_{ab}.
\]
We write
\[
\mathfrakh
:=
\{H\in\C^{3\times3}:H^\dagger=H,\ \tr H=0\}
\]
for the Hermitian traceless generator space.
The Lie algebra $\fraksu$ is identified with $i\mathfrakh$ in the usual
physics convention.
\end{definition}

\begin{definition}[Realized node colors and local edge data]
\label{def:realized-node-color}
Fix the color projection
\[
\Pi_{\mathrm{col}}:\R^4\to\R^3
\]
and the velocity squashing map
\[
\psi(v):=\frac{V_{\mathrm{alg}}}{V_{\mathrm{alg}}+\|v\|}\,v.
\]
For each realized node \(n\in\mcV^{(N,T)}(\omega)\), define its raw and
normalized local color vectors by
\[
\tilde{\mathbf c}_n(\omega):=\Pi_{\mathrm{col}}\psi(v_n(\omega))\in\R^3\subset\C^3,
\]
\[
\mathbf c_n(\omega)
:=
\frac{\tilde{\mathbf c}_n(\omega)}
{\sqrt{\tilde{\mathbf c}_n(\omega)^\dagger\tilde{\mathbf c}_n(\omega)+\epsilon_c^2}}
\in\C^3.
\]
For each IG edge \(e=(n\to m)\), let \(\theta_e(\omega)\in\R\) denote the
realized cloning-score phase extracted from the run history, and define its
Lie-direction coordinates by
\[
\hat n_e^{(a)}(\omega)
:=
\frac{
\Re\!\bigl(\mathbf c_m(\omega)^\dagger\lambda_a\,\mathbf c_n(\omega)\bigr)}
{\Bigl(\sum_{b=1}^8
\Re\!\bigl(\mathbf c_m(\omega)^\dagger\lambda_b\,\mathbf c_n(\omega)\bigr)^2
\Bigr)^{1/2}+\epsilon_c},
\qquad
1\le a\le 8.
\]
For each IA edge \(e\), let \(R_e^{\mathrm{IA}}(\omega)\in SU(3)\) be the
stored non-abelian attribution rotation carried by that edge.
\end{definition}

\begin{definition}[Realized edge transports]
\label{def:realized-edge-transports}
For each oriented edge \(e\) of the realized Fractal regulator, define its
\(SU(3)\) transport by the following rules.
\begin{enumerate}[leftmargin=1.5em,label=\textnormal{(\roman*)}]
\item If \(e\in\mcE_{\mathrm{CST}}^{(N,T)}(\omega)\), set
\[
U_e(\omega):=I_3.
\]
\item If \(e\in\mcE_{\mathrm{IG}}^{(N,T)}(\omega)\), define the Hermitian
generator
\[
\Theta_e(\omega)
:=
\theta_e(\omega)\sum_{a=1}^8\hat n_e^{(a)}(\omega)\lambda_a\in\mathfrakh,
\]
and set
\[
U_e(\omega):=\exp\!\bigl(i\Theta_e(\omega)\bigr)\in SU(3).
\]
\item If \(e\in\mcE_{\mathrm{IA}}^{(N,T)}(\omega)\), set
\[
U_e(\omega):=R_e^{\mathrm{IA}}(\omega)\in SU(3).
\]
\item For the reversed orientation, define
\[
U_{e^{-1}}(\omega):=\bigl(U_e(\omega)\bigr)^{-1}.
\]
\end{enumerate}
\end{definition}

\begin{lemma}[Realized link transport lies in \(SU(3)\)]
\label{lem:realized-link-su3}
For every \(\omega\in\Omega_{N,T}\) and every oriented edge \(e\) of the
realized Fractal regulator, \(U_e(\omega)\in SU(3)\).
\end{lemma}

\begin{proof}
For CST edges this is immediate because \(U_e(\omega)=I_3\). For IG edges,
\(\Theta_e(\omega)\in\mathfrakh\), so \(i\Theta_e(\omega)\) is skew-Hermitian
and traceless; therefore \(\exp(i\Theta_e(\omega))\in SU(3)\). For IA edges,
membership in \(SU(3)\) is part of the stored attribution payload.
\end{proof}

\begin{definition}[Wilson observables on the realized Fractal regulator]
\label{def:realized-wilson-observables}
Let \(\omega\in\Omega_{N,T}\).
\begin{enumerate}[leftmargin=1.5em,label=\textnormal{(\roman*)}]
\item For an oriented loop \(\gamma=(e_1,\dots,e_m)\) in the realized regulator,
define the Wilson loop transport and normalized Wilson observable by
\[
U_\gamma(\omega):=\prod_{r=1}^m U_{e_r}(\omega),
\qquad
W_\gamma(\omega):=\frac13\tr U_\gamma(\omega).
\]
\item For a plaquette \(P\in\mcP_{\mathrm{FG}}^{(N,T)}(\omega)\), define the
plaquette holonomy
\[
U_P(\omega):=\prod_{e\in\partial P}U_e(\omega).
\]
Whenever \(A_P(\omega)>0\) and \(U_P(\omega)\) lies in the principal logarithm
domain, define the discrete field strength by
\[
F_P(\omega):=\frac{1}{ig_3A_P(\omega)}\log U_P(\omega)\in\mathfrakh.
\]
\item The finite Wilson action is the random variable
\[
S_{\mathrm W}^{(N,T)}(\omega)
:=
\frac{6}{g_3^2}
\sum_{P\in\mcP_{\mathrm{FG}}^{(N,T)}(\omega)}
\left(
1-\frac13\Re\tr U_P(\omega)
\right).
\]
\end{enumerate}
\end{definition}

\begin{definition}[Gauge maps on the realized regulator]
\label{def:realized-gauge-map}
A gauge map on the realized regulator is a function
\[
\Omega:\mcV^{(N,T)}(\omega)\to SU(3).
\]
Its action on a link \(e=(n\to m)\) is
\[
U_e^\Omega(\omega):=
\Omega(m)\,U_e(\omega)\,\Omega(n)^{-1}.
\]
\end{definition}

\begin{proposition}[Bridge from the Hermitian convention to the standard
\(\mathfrak{su}(3)\) convention]
\label{prop:su3-convention-bridge}
The continuum and lattice objects of this paper are exactly the standard
\(SU(3)\) Yang--Mills objects after the usual identification
\[
\fraksu=i\mathfrakh.
\]
More precisely:
\begin{enumerate}[leftmargin=1.5em,label=\textnormal{(\roman*)}]
\item For every differentiable color field \(c:\R^4\to\C^3\), define the
anti-Hermitian connection and curvature by
\[
\mathcal A_\mu[c]:=iA_\mu[c]\in\fraksu,
\qquad
\mathcal F_{\mu\nu}[c]:=iF_{\mu\nu}[c]\in\fraksu.
\]
Then
\[
\mathcal F_{\mu\nu}[c]
=
\partial_\mu\mathcal A_\nu[c]-\partial_\nu\mathcal A_\mu[c]
+g_3[\mathcal A_\mu[c],\mathcal A_\nu[c]],
\]
which is the standard anti-Hermitian \(SU(3)\) curvature formula.
\item For the stationary continuum field \(F^\infty\) of this paper, the local
and global Yang--Mills densities agree with the standard anti-Hermitian
normalization:
\[
\tr\!\bigl((F_{\mu\nu}^\infty(x))^2\bigr)
=
-\tr\!\bigl((\mathcal F_{\mu\nu}^\infty(x))^2\bigr),
\qquad
\mathcal F_{\mu\nu}^\infty:=iF_{\mu\nu}^\infty,
\]
so the continuum comparison density used in this paper is exactly the usual
positive \(SU(3)\) Yang--Mills density written in the anti-Hermitian
convention.
\item For every realized edge and plaquette, if
\[
\mathcal A_e:=i\Theta_e\in\fraksu
\quad\text{on IG edges},\qquad
\mathcal F_P:=iF_P\in\fraksu,
\]
then
\[
U_e=\exp(\mathcal A_e)
\quad\text{on IG edges},\qquad
U_P=\exp\!\bigl(g_3A_P\mathcal F_P\bigr)
\]
whenever the principal logarithm is defined, so the group-valued Wilson
transport variables are unchanged by passing from the Hermitian convention to
the standard anti-Hermitian \(\fraksu\) convention.
\end{enumerate}
\end{proposition}

\begin{proof}
For \textnormal{(i)}, \(\mathcal A_\mu[c]=iA_\mu[c]\) is skew-Hermitian and
traceless because \(A_\mu[c]\in\mathfrakh\); the same holds for
\(\mathcal F_{\mu\nu}[c]=iF_{\mu\nu}[c]\). Using
\[
[\mathcal A_\mu[c],\mathcal A_\nu[c]]
=[iA_\mu[c],iA_\nu[c]]
=-[A_\mu[c],A_\nu[c]],
\]
one computes
\begin{align*}
\partial_\mu\mathcal A_\nu[c]-\partial_\nu\mathcal A_\mu[c]
+g_3[\mathcal A_\mu[c],\mathcal A_\nu[c]]
&=
i\partial_\mu A_\nu[c]-i\partial_\nu A_\mu[c]-g_3[A_\mu[c],A_\nu[c]]\\
&=
i\Bigl(\partial_\mu A_\nu[c]-\partial_\nu A_\mu[c]
+ig_3[A_\mu[c],A_\nu[c]]\Bigr)\\
&=
iF_{\mu\nu}[c]
=
\mathcal F_{\mu\nu}[c].
\end{align*}

For \textnormal{(ii)}, \(\mathcal F_{\mu\nu}^\infty=iF_{\mu\nu}^\infty\) gives
\[
(\mathcal F_{\mu\nu}^\infty)^2
=
- (F_{\mu\nu}^\infty)^2,
\]
hence
\[
-\tr\!\bigl((\mathcal F_{\mu\nu}^\infty)^2\bigr)
=
\tr\!\bigl((F_{\mu\nu}^\infty)^2\bigr).
\]
Therefore the localized continuum comparison functional built from
\(\tr((F_{\mu\nu}^\infty)^2)\) is exactly the usual positive Yang--Mills
density in anti-Hermitian \(\fraksu\) variables.

For \textnormal{(iii)}, the IG-edge case is immediate from
Definition~\ref{def:realized-edge-transports}. The plaquette statement is
exactly the logarithmic identity proved below in
Proposition~\ref{prop:random-plaquette-wilson-expansion}.
\end{proof}

\begin{proposition}[Exact plaquette exponential identity and Wilson expansion]
\label{prop:random-plaquette-wilson-expansion}
Fix \(\omega\in\Omega_{N,T}\), and let
\(P\in\mcP_{\mathrm{FG}}^{(N,T)}(\omega)\) be such that \(A_P(\omega)>0\) and
\(U_P(\omega)\) lies in the principal logarithm domain.
Then
\[
U_P(\omega)=\exp\!\bigl(i g_3 A_P(\omega)F_P(\omega)\bigr).
\]
If in addition
\[
\|g_3A_P(\omega)F_P(\omega)\|_{\mathrm{op}}\le \delta
\]
for some fixed $\delta>0$, then there exists a constant $C_\delta<\infty$,
depending only on $\delta$ and the matrix norm convention, such that
\[
1-\frac13\Re\tr U_P(\omega)
=
\frac{g_3^2A_P(\omega)^2}{6}\,
\tr\!\bigl((F_P(\omega))^2\bigr)
+R_P(\omega),
\]
with
\[
|R_P(\omega)|
\le
C_\delta\,g_3^3A_P(\omega)^3\|F_P(\omega)\|_{\mathrm{op}}^3.
\]
\end{proposition}

\begin{proof}
By Definition~\ref{def:realized-wilson-observables},
\[
F_P(\omega)
=
\frac{1}{ig_3A_P(\omega)}\log U_P(\omega),
\]
so multiplication by \(ig_3A_P(\omega)\) gives
\[
\log U_P(\omega)=ig_3A_P(\omega)F_P(\omega).
\]
Exponentiating both sides yields the exact identity
\[
U_P(\omega)=\exp\!\bigl(i g_3 A_P(\omega)F_P(\omega)\bigr).
\]

Set
\[
X_P(\omega):=ig_3A_P(\omega)F_P(\omega).
\]
Because \(F_P(\omega)\in\mathfrakh\), the matrix \(X_P(\omega)\) is skew-Hermitian and
traceless.
The exponential series gives
\[
e^{X_P(\omega)}
=
I+X_P(\omega)+\frac12(X_P(\omega))^2+\mathcal R_P(\omega),
\]
where the third-order remainder satisfies the standard matrix bound
\[
\|\mathcal R_P(\omega)\|_{\mathrm{op}}
\le
C_\delta\|X_P(\omega)\|_{\mathrm{op}}^3
\]
whenever \(\|X_P(\omega)\|_{\mathrm{op}}\le\delta\).
Taking traces and using \(\tr X_P(\omega)=0\), we obtain
\[
\tr U_P(\omega)
=
3+\frac12\tr\!\bigl((X_P(\omega))^2\bigr)+\tr\mathcal R_P(\omega).
\]
Since \((ig_3A_P(\omega))^2=-g_3^2A_P(\omega)^2\),
\[
\tr\!\bigl((X_P(\omega))^2\bigr)
=
-g_3^2A_P(\omega)^2\tr\!\bigl((F_P(\omega))^2\bigr).
\]
Therefore
\[
\Re\tr U_P(\omega)
=
3-\frac{g_3^2A_P(\omega)^2}{2}\tr\!\bigl((F_P(\omega))^2\bigr)
+\Re\tr\mathcal R_P(\omega).
\]
Subtracting from $1$ after division by $3$ gives
\[
1-\frac13\Re\tr U_P(\omega)
=
\frac{g_3^2A_P(\omega)^2}{6}\tr\!\bigl((F_P(\omega))^2\bigr)
-\frac13\Re\tr\mathcal R_P(\omega).
\]
It remains to estimate the last term.
In fixed dimension, $|\tr Y|\le 3\|Y\|_{\mathrm{op}}$ for every
$3\times3$ matrix $Y$, so
\[
\left|-\frac13\Re\tr\mathcal R_P(\omega)\right|
\le
\|\mathcal R_P(\omega)\|_{\mathrm{op}}
\le
C_\delta\|X_P(\omega)\|_{\mathrm{op}}^3.
\]
Since \(\|X_P(\omega)\|_{\mathrm{op}}=g_3A_P(\omega)\|F_P(\omega)\|_{\mathrm{op}}\),
the claimed bound follows.
\end{proof}

\begin{proposition}[Gauge invariance of realized plaquette traces]
\label{prop:realized-gauge-invariance}
For every realized plaquette \(P\in\mcP_{\mathrm{FG}}^{(N,T)}(\omega)\) and
every gauge map \(\Omega\) on the realized regulator,
\[
\tr U_P^\Omega(\omega)=\tr U_P(\omega).
\]
Consequently,
\[
S_{\mathrm W}^{(N,T),\Omega}(\omega)=S_{\mathrm W}^{(N,T)}(\omega).
\]
\end{proposition}

\begin{proof}
Write the plaquette boundary as
\[
v_0\to v_1\to v_2\to v_3\to v_4=v_0.
\]
Under the gauge action from Definition~\ref{def:realized-gauge-map},
\[
U_P^\Omega(\omega)
=
\prod_{k=1}^4
\Omega(v_k)\,U_{(v_{k-1}\to v_k)}(\omega)\,\Omega(v_{k-1})^{-1}.
\]
All consecutive inner factors telescope, leaving
\[
U_P^\Omega(\omega)
=
\Omega(v_0)\,U_P(\omega)\,\Omega(v_0)^{-1}.
\]
Trace cyclicity gives the first claim, and summing over all plaquettes gives
the Wilson-action invariance.
\end{proof}

\paragraph{Stage I purpose.}
The purpose of the next package is to upgrade the finite random gauge
construction from a measurable pushforward law on realized gauge data into a
fibered finite-volume gauge measure theory with a precise configuration-space
theorem, disintegration over realized regulators, a gauge-invariance theorem
for the conditional laws, an exact effective-action representation under an
explicit absolute-continuity hypothesis, and a Wilson-comparison theorem under
an explicit remainder/oscillation hypothesis. This is the measure-theoretic
foundation needed before one can honestly build a nonabelian Euclidean gauge
theory on the Fractal-Gas regulator.

\paragraph{Role of Stage I inside the revised paper.}
The realized Fractal regulator, realized edge transports, plaquette holonomies,
Wilson action, and pushforward finite gauge law introduced above are the right
first move, but they are not yet enough to support Yang--Mills QFT. The precise
gap is that the induced gauge law is not yet written as a gauge measure
relative to a standard reference measure, no exact effective action has been
defined, no rigorous comparison theorem has been proved between the induced law
and the Wilson law, and no fiberwise gauge-invariance statement has been
isolated at the level of the conditional laws. The next Stage I package fixes
exactly those four problems.

\begin{definition}[Fixed finite horizon and basic gauge object]
\label{def:fixed-finite-horizon}
Fix once and for all a swarm size \(N\), a recorded horizon \(T\), and the
compact gauge group \(G:=SU(3)\). The swarm-history law is
\[
(\Omega_{N,T},\mathcal F_{N,T},\mathbb P_{N,T}).
\]
Each history \(\omega\in\Omega_{N,T}\) determines a realized Fractal regulator,
written
\[
R(\omega)\in\mathscr R_{N,T},
\]
together with stored positive-edge transports
\[
\{U_e(\omega)\}_{e\in E^+(R(\omega))}\in G^{E^+(R(\omega))}.
\]
The reversed-edge transport is always defined by \(U_{\bar e}(\omega)=U_e(\omega)^{-1}\).
\end{definition}

\begin{remark}[Basic regulator data]
\label{rem:basic-regulator-data}
The regulator data \(R(\omega)\) records the finite vertex set, the stored
positive oriented edges, the oriented plaquettes, the incidence maps, the
plaquette basepoints \(x_P\), bivectors \(\Sigma_P\), areas \(A_P\), edge and
plaquette types, and any additional bounded local combinatorial or geometric
payload needed to evaluate the finite Wilson action already defined above.
\end{remark}

\begin{definition}[Realized regulator-data space]
\label{def:realized-regulator-data-space}
Let \(M_V(N,T)\), \(M_E(N,T)\), and \(M_P(N,T)\) be deterministic upper bounds
on the numbers of realized vertices, stored positive edges, and plaquettes that
can occur at size \(N\) and horizon \(T\). For each triple
\((m_V,m_E,m_P)\), let \(\mathscr C_{N,T}(m_V,m_E,m_P)\) denote the finite
product of Euclidean factors, compact intervals, and finite label sets needed
to store the bounded auxiliary local payload from
Remark~\ref{rem:basic-regulator-data}. Define
\[
\mathscr R_{N,T}
:=
\bigsqcup_{m_V=0}^{M_V(N,T)}
\bigsqcup_{m_E=0}^{M_E(N,T)}
\bigsqcup_{m_P=0}^{M_P(N,T)}
\mathscr R_{N,T}(m_V,m_E,m_P),
\]
where
\begin{align*}
\mathscr R_{N,T}(m_V,m_E,m_P)
:={}&
(\R\times\R^4)^{m_V}
\times \{0,1,2\}^{m_E}
\times \{1,\dots,m_V\}^{2m_E}
\times \{1,\dots,m_E\}^{4m_P}\\
&\times \{-1,+1\}^{4m_P}
\times (\R^4)^{m_P}
\times (\bigwedge^2\R^4)^{m_P}
\times [0,\infty)^{m_P}
\times \mathscr C_{N,T}(m_V,m_E,m_P).
\end{align*}
Equip \(\mathscr R_{N,T}\) with the disjoint-union product \(\sigma\)-algebra.
\end{definition}

\begin{proposition}[Standard Borelness of the regulator space]
\label{prop:regulator-space-standard-borel}
The space \(\mathscr R_{N,T}\) is a standard Borel space.
\end{proposition}

\begin{proof}
For each fixed \((m_V,m_E,m_P)\), the component
\(\mathscr R_{N,T}(m_V,m_E,m_P)\) is a finite product of Euclidean spaces,
compact intervals, and finite sets, hence is standard Borel. The ranges of
\(m_V,m_E,m_P\) are finite, so the full disjoint union is again standard
Borel.
\end{proof}

\begin{definition}[Fiber gauge configuration space]
\label{def:fiber-gauge-configuration-space}
For \(R\in\mathscr R_{N,T}(m_V,m_E,m_P)\), let \(E^+(R)\) be its stored
positive-edge set. The corresponding gauge fiber is
\[
\mathscr U(R):=SU(3)^{E^+(R)}\cong SU(3)^{m_E}.
\]
An element \(U\in\mathscr U(R)\) is a family \(\{U_e\}_{e\in E^+(R)}\) of
stored positive-edge transports.
\end{definition}

\begin{proposition}[Compact Polish structure of the fibers]
\label{prop:fiber-compact-polish}
For every \(R\in\mathscr R_{N,T}\), the fiber \(\mathscr U(R)\) is a compact
metrizable group and hence a compact Polish space.
\end{proposition}

\begin{proof}
The compact Lie group \(SU(3)\) is compact and metrizable. Every fiber is a
finite Cartesian product of copies of \(SU(3)\).
\end{proof}

\begin{definition}[Fiberwise Wilson action]
\label{def:fiberwise-wilson-action}
For \(R\in\mathscr R_{N,T}\) and \(U\in\mathscr U(R)\), write
\[
U_\gamma(U):=\prod_{e\in\gamma} U_e^{\epsilon(\gamma,e)}
\]
for the ordered holonomy of a finite loop \(\gamma\), and define the normalized
Wilson loop observable by
\[
W_\gamma(R,U):=\frac13\tr U_\gamma(U).
\]
Also write
\[
U_P(U):=\prod_{e\in\partial P} U_e^{\epsilon(P,e)}
\]
for the ordered plaquette holonomy, where \(\epsilon(P,e)\in\{-1,+1\}\) records
whether the stored orientation of \(e\) agrees with the boundary orientation of
\(P\). Define
\[
S_{\mathrm W}(R,U)
:=
\frac{6}{g_3^2}
\sum_{P\in P(R)}
\left(
1-\frac13\Re\tr U_P(U)
\right).
\]
For a realized history \(\omega\), this coincides with
\(S_{\mathrm W}^{(N,T)}(\omega)\) from
Definition~\ref{def:realized-wilson-observables}.
\end{definition}

\begin{proposition}[Measurability and boundedness of Wilson observables]
\label{prop:wilson-observable-measurable-bounded}
For every fixed finite loop \(\gamma\) and plaquette \(P\), the maps
\[
(R,U)\mapsto W_\gamma(R,U),\qquad
(R,U)\mapsto \Re\tr U_P(U),\qquad
(R,U)\mapsto S_{\mathrm W}(R,U)
\]
are measurable on the total gauge space defined below. Moreover, all Wilson
loops satisfy
\[
|W_\gamma(R,U)|\le 1.
\]
\end{proposition}

\begin{proof}
Each Wilson loop or plaquette trace is obtained from finitely many coordinate
projections on \(SU(3)\) using matrix inversion, multiplication, and trace, all
of which are continuous. The Wilson action is a finite sum of measurable
plaquette terms. Since \(SU(3)\subset U(3)\), every loop transport has operator
norm \(1\), so its trace is uniformly bounded.
\end{proof}

\begin{definition}[Total fibered gauge space]
\label{def:total-fibered-gauge-space}
Define
\[
\mathscr X_{N,T}
:=
\bigsqcup_{m_V=0}^{M_V(N,T)}
\bigsqcup_{m_E=0}^{M_E(N,T)}
\bigsqcup_{m_P=0}^{M_P(N,T)}
\Bigl(\mathscr R_{N,T}(m_V,m_E,m_P)\times SU(3)^{m_E}\Bigr).
\]
Equivalently,
\[
\mathscr X_{N,T}=\{(R,U):R\in \mathscr R_{N,T},\ U\in\mathscr U(R)\},
\]
equipped with the corresponding disjoint-union product \(\sigma\)-algebra.
\end{definition}

\begin{proposition}[Standard Borelness of the total gauge space]
\label{prop:total-gauge-space-standard-borel}
The total gauge space \(\mathscr X_{N,T}\) is a standard Borel space.
\end{proposition}

\begin{proof}
Each component
\(\mathscr R_{N,T}(m_V,m_E,m_P)\times SU(3)^{m_E}\) is a product of standard
Borel spaces, hence standard Borel. The full space is a finite disjoint union
of such components.
\end{proof}

\begin{definition}[Regulator extraction and gauge construction]
\label{def:regulator-extraction-gauge-construction}
Let
\[
\mathfrak R^{(N,T)}:\Omega_{N,T}\to \mathscr R_{N,T}
\]
be the measurable map extracting the realized Fractal regulator data. Let
\[
\mathfrak U^{(N,T)}:\Omega_{N,T}\to \bigsqcup_{R\in\mathscr R_{N,T}}\mathscr U(R)
\]
be the measurable map assigning the stored positive-edge transports from
Definition~\ref{def:realized-edge-transports}. Define the combined
gauge-construction map
\[
\mathfrak G^{(N,T)}(\omega)
:=
\bigl(\mathfrak R^{(N,T)}(\omega),\mathfrak U^{(N,T)}(\omega)\bigr)\in \mathscr X_{N,T}.
\]
\end{definition}

\begin{theorem}[Induced finite-volume gauge law]
\label{thm:induced-finite-volume-gauge-law}
The pushforward measure
\[
\mu_{N,T}:=\bigl(\mathfrak G^{(N,T)}\bigr)_{\#}\mathbb P_{N,T}\in\mathcal P(\mathscr X_{N,T})
\]
is well defined. For every bounded measurable observable
\(\Phi:\mathscr X_{N,T}\to\R\),
\[
\int_{\mathscr X_{N,T}}\Phi(R,U)\,\mu_{N,T}(dR,dU)
=
\E_{\mathbb P_{N,T}}\bigl[\Phi(\mathfrak G^{(N,T)}(\omega))\bigr].
\]
\end{theorem}

\begin{proof}
Proposition~\ref{prop:fractal-regulator-measurable} gives measurability of the
realized-regulator map. Definition~\ref{def:realized-edge-transports} builds
the stored edge transports from finitely many measurable algebraic operations
on the recorded history, so \(\mathfrak U^{(N,T)}\) is measurable. Therefore
\(\mathfrak G^{(N,T)}\) is measurable as a map into the standard Borel space
\(\mathscr X_{N,T}\), and the pushforward measure is well defined. The integral
identity is the definition of pushforward.
\end{proof}

\begin{corollary}[Existence of all finite Wilson moments]
\label{cor:finite-wilson-moments}
Every bounded gauge observable generated by finitely many Wilson loops and
plaquette traces is \(\mu_{N,T}\)-integrable. In particular, all finite joint
moments of Wilson loops exist.
\end{corollary}

\begin{proof}
By Proposition~\ref{prop:wilson-observable-measurable-bounded}, such
observables are bounded and measurable on \(\mathscr X_{N,T}\).
\end{proof}

\begin{theorem}[Disintegration over realized regulators]
\label{thm:disintegration-over-regulators}
Let
\[
\lambda_{N,T}:=\bigl(\mathfrak R^{(N,T)}\bigr)_{\#}\mathbb P_{N,T}\in\mathcal P(\mathscr R_{N,T}).
\]
Then there exists a \(\lambda_{N,T}\)-almost everywhere uniquely determined
measurable family of probability measures
\[
R\mapsto \mu_{N,T}^R\in\mathcal P(\mathscr U(R))
\]
such that for every bounded measurable \(\Phi\) on \(\mathscr X_{N,T}\),
\[
\int_{\mathscr X_{N,T}}\Phi(R,U)\,\mu_{N,T}(dR,dU)
=
\int_{\mathscr R_{N,T}}
\left(
\int_{\mathscr U(R)}\Phi(R,U)\,\mu_{N,T}^R(dU)
\right)
\lambda_{N,T}(dR).
\]
\end{theorem}

\begin{proof}
The projection \((R,U)\mapsto R\) is measurable between standard Borel spaces.
Rokhlin disintegration therefore applies to \(\mu_{N,T}\), yielding a
measurable kernel \(R\mapsto\mu_{N,T}^R\) supported on the fiber \(\mathscr U(R)\).
\end{proof}

\begin{remark}[What has been proved unconditionally]
\label{rem:unconditional-part}
Up to Theorem~\ref{thm:disintegration-over-regulators}, everything is
unconditional. The only inputs are finite horizon, finite swarm size,
measurability of the regulator and gauge-construction maps, and compactness of
the gauge group. No Gibbs structure, Yang--Mills comparison, or absolute
continuity has been assumed yet.
\end{remark}

\begin{definition}[Fiber gauge group]
\label{def:fiber-gauge-group}
For each realized regulator \(R\), define
\[
\mathcal G(R):=SU(3)^{V(R)}.
\]
For \(g\in\mathcal G(R)\) and \(U\in\mathscr U(R)\), define
\[
(g\cdot U)_e:=g_{t(e)}U_e g_{s(e)}^{-1},
\]
where \(s(e)\) and \(t(e)\) are the source and target vertices of the stored
oriented edge \(e\).
\end{definition}

\begin{proposition}[Fiberwise gauge action]
\label{prop:fiberwise-gauge-action}
For each fixed \(R\), the map
\[
\mathcal G(R)\times \mathscr U(R)\to \mathscr U(R),\qquad (g,U)\mapsto g\cdot U
\]
is continuous.
\end{proposition}

\begin{proof}
It is a finite product of multiplication and inversion maps in the compact Lie
group \(SU(3)\).
\end{proof}

\begin{definition}[Gauge-invariant observable]
\label{def:gauge-invariant-observable}
A measurable function \(\Phi\) on \(\mathscr X_{N,T}\) is called
gauge-invariant if for every \(R\in\mathscr R_{N,T}\), every
\(g\in\mathcal G(R)\), and every \(U\in\mathscr U(R)\),
\[
\Phi(R,g\cdot U)=\Phi(R,U).
\]
Wilson loops and the Wilson action are gauge-invariant in this sense.
\end{definition}

\begin{assumption}[Gauge-equivariant lift on the swarm history space]
\label{ass:gauge-equivariant-lift}
For each realized regulator \(R\) and each local gauge map
\(g\in\mathcal G(R)\), there exists a measurable transformation
\[
\widetilde g:\Omega_{N,T}\to\Omega_{N,T}
\]
such that:
\begin{enumerate}[leftmargin=1.5em,label=\textnormal{(\roman*)}]
\item \textbf{Regulator invariance:}
\[
\mathfrak R^{(N,T)}(\widetilde g\omega)=\mathfrak R^{(N,T)}(\omega).
\]
\item \textbf{Gauge-equivariance of the edge assignment:}
\[
\mathfrak U^{(N,T)}(\widetilde g\omega)=g\cdot \mathfrak U^{(N,T)}(\omega).
\]
\item \textbf{Invariance of the swarm law:}
\[
\mathbb P_{N,T}\circ \widetilde g^{-1}=\mathbb P_{N,T}.
\]
\end{enumerate}
\end{assumption}

\begin{theorem}[Fiberwise gauge invariance of the induced gauge law]
\label{thm:fiberwise-gauge-invariance}
Under Assumption~\ref{ass:gauge-equivariant-lift}, for
\(\lambda_{N,T}\)-almost every realized regulator \(R\), the conditional law
\(\mu_{N,T}^R\) is gauge-invariant:
\[
\mu_{N,T}^R(A)=\mu_{N,T}^R(g\cdot A)
\]
for every Borel set \(A\subseteq\mathscr U(R)\) and every \(g\in\mathcal G(R)\).
Equivalently, for every bounded measurable \(f:\mathscr U(R)\to\R\),
\[
\int_{\mathscr U(R)}f(U)\,\mu_{N,T}^R(dU)
=
\int_{\mathscr U(R)}f(g\cdot U)\,\mu_{N,T}^R(dU).
\]
\end{theorem}

\begin{proof}
Fix a realized regulator \(R\) and a bounded measurable
\(f:\mathscr U(R)\to\R\). Define the corresponding observable on
\(\Omega_{N,T}\) by
\[
F_f(\omega):=
f\bigl(\mathfrak U^{(N,T)}(\omega)\bigr)\mathbf 1_{\{\mathfrak R^{(N,T)}(\omega)=R\}}.
\]
By Assumption~\ref{ass:gauge-equivariant-lift},
\[
F_f(\widetilde g\omega)
=
f\bigl(g\cdot \mathfrak U^{(N,T)}(\omega)\bigr)
\mathbf 1_{\{\mathfrak R^{(N,T)}(\omega)=R\}}.
\]
Since \(\mathbb P_{N,T}\) is invariant under \(\widetilde g\),
\[
\E[F_f]=\E[F_f\circ \widetilde g].
\]
Now fix a component \(\mathscr R_{N,T}(m_V,m_E,m_P)\). Because the vertex labels
in that component are identified with \(\{1,\dots,m_V\}\), the same
\(g\in SU(3)^{m_V}\cong\mathcal G(R)\) acts on every fiber over that component.
For an arbitrary measurable set \(B\subseteq\mathscr R_{N,T}(m_V,m_E,m_P)\),
define
\[
\Phi_{f,B}(R',U')
:=
\mathbf 1_B(R')\,f(U').
\]
Applying the previous expectation identity with \(\Phi_{f,B}\) and then using
Theorem~\ref{thm:disintegration-over-regulators} gives
\[
\int_B
\left[
\int_{\mathscr U(R')} f(U')\,\mu_{N,T}^{R'}(dU')
-
\int_{\mathscr U(R')} f(g\cdot U')\,\mu_{N,T}^{R'}(dU')
\right]
\lambda_{N,T}(dR')
=0.
\]
Since \(B\) is arbitrary, the bracketed integrand vanishes for
\(\lambda_{N,T}\)-almost every \(R'\) in that component. Summing over the
finitely many components yields
\[
\int_{\mathscr U(R)}f(U)\,\mu_{N,T}^R(dU)
=
\int_{\mathscr U(R)}f(g\cdot U)\,\mu_{N,T}^R(dU)
\]
for \(\lambda_{N,T}\)-almost every \(R\). This is exactly gauge invariance of
the conditional law.
\end{proof}

\begin{corollary}[Gauge-invariant observables descend to the quotient]
\label{cor:gauge-invariant-observables-descend}
Under Assumption~\ref{ass:gauge-equivariant-lift}, expectations of
gauge-invariant observables depend only on the gauge orbit of the configuration
in each fiber.
\end{corollary}

\begin{proof}
If \(\Phi(R,\cdot)\) is gauge-invariant, then
\(\Phi(R,g\cdot U)=\Phi(R,U)\) pointwise, and the claim follows immediately
from Theorem~\ref{thm:fiberwise-gauge-invariance}.
\end{proof}

\begin{definition}[Fiberwise reference Haar measure]
\label{def:fiberwise-reference-haar}
For each realized regulator \(R\), let
\[
\nu_R:=\bigotimes_{e\in E^+(R)} d\mathrm{Haar}(U_e)\in\mathcal P(\mathscr U(R))
\]
be the product Haar measure on the fiber \(\mathscr U(R)\).
\end{definition}

\begin{proposition}[Gauge invariance of the fiber Haar measures]
\label{prop:gauge-invariance-fiber-haar}
For every realized regulator \(R\) and every \(g\in\mathcal G(R)\),
\[
g_{\#}\nu_R=\nu_R.
\]
\end{proposition}

\begin{proof}
On each edge variable, left multiplication at the target and right
multiplication by the inverse at the source preserve Haar measure. Finite
products preserve the invariance.
\end{proof}

\begin{proposition}[Measurable reference-measure kernel]
\label{prop:measurable-reference-kernel}
The assignment
\[
R\mapsto \nu_R
\]
is a measurable probability kernel from \(\mathscr R_{N,T}\) to
\(\mathscr X_{N,T}\).
\end{proposition}

\begin{proof}
On each component \(\mathscr R_{N,T}(m_V,m_E,m_P)\), the measure \(\nu_R\) is
just the product Haar measure on \(SU(3)^{m_E}\), independent of the remaining
regulator coordinates. This gives a measurable kernel on each component, and
hence on the full finite disjoint union.
\end{proof}

\begin{assumption}[Fiberwise absolute continuity]
\label{ass:fiberwise-absolute-continuity}
For \(\lambda_{N,T}\)-almost every regulator \(R\), the conditional induced
gauge law \(\mu_{N,T}^R\) is absolutely continuous with respect to the fiber
Haar measure \(\nu_R\):
\[
\mu_{N,T}^R\ll \nu_R.
\]
\end{assumption}

\begin{theorem}[Exact effective-action representation]
\label{thm:exact-effective-action}
Under Assumption~\ref{ass:fiberwise-absolute-continuity}, there exists a
jointly measurable nonnegative density
\[
\rho_{N,T}(R,U)=\frac{d\mu_{N,T}^R}{d\nu_R}(U)
\]
such that
\[
\mu_{N,T}(dR,dU)=\lambda_{N,T}(dR)\,\rho_{N,T}(R,U)\,\nu_R(dU).
\]
Define
\[
S_{\mathrm{eff}}^{(N,T)}(R,U):=
\begin{cases}
-\log \rho_{N,T}(R,U), & \rho_{N,T}(R,U)>0,\\[4pt]
+\infty, & \rho_{N,T}(R,U)=0.
\end{cases}
\]
Then
\[
\mu_{N,T}(dR,dU)=\lambda_{N,T}(dR)\,
e^{-S_{\mathrm{eff}}^{(N,T)}(R,U)}\,\nu_R(dU).
\]
\end{theorem}

\begin{proof}
Theorem~\ref{thm:disintegration-over-regulators} gives the measurable kernel
\(R\mapsto \mu_{N,T}^R\), and
Proposition~\ref{prop:measurable-reference-kernel} gives the measurable kernel
\(R\mapsto \nu_R\). Assumption~\ref{ass:fiberwise-absolute-continuity} allows
the Radon--Nikodym theorem for measurable kernels to produce a jointly
measurable density \(\rho_{N,T}\). The exponential form is then just the
definition of \(S_{\mathrm{eff}}^{(N,T)}\).
\end{proof}

\begin{remark}[Scope of the effective action]
\label{rem:effective-action-scope}
Theorem~\ref{thm:exact-effective-action} upgrades the pushforward law to a
finite-volume gauge measure with an exact effective action. It does not yet
identify \(S_{\mathrm{eff}}^{(N,T)}\) with the Wilson action.
\end{remark}

\begin{definition}[Fiberwise Wilson measure]
\label{def:fiberwise-wilson-measure}
For each realized regulator \(R\), define the partition function
\[
Z_W(R):=\int_{\mathscr U(R)} e^{-S_{\mathrm W}(R,U)}\,\nu_R(dU).
\]
Since \(\mathscr U(R)\) is compact and \(S_{\mathrm W}(R,\cdot)\) is bounded
below and measurable, \(0<Z_W(R)<\infty\). Define the Wilson fiber measure
\[
\mu_W^R(dU):=Z_W(R)^{-1}e^{-S_{\mathrm W}(R,U)}\,\nu_R(dU).
\]
Define the mixed Wilson law on the total space by
\[
\mu_W^{(N,T)}(dR,dU):=\lambda_{N,T}(dR)\,\mu_W^R(dU).
\]
\end{definition}

\begin{proposition}[Gauge invariance of the Wilson fibers]
\label{prop:gauge-invariance-wilson-fibers}
For every realized regulator \(R\), the Wilson measure \(\mu_W^R\) is
gauge-invariant.
\end{proposition}

\begin{proof}
The Haar measure \(\nu_R\) is gauge-invariant by
Proposition~\ref{prop:gauge-invariance-fiber-haar}. For fixed \(R\), the same
plaquette-boundary telescoping argument used in
Proposition~\ref{prop:realized-gauge-invariance} shows that
\(S_{\mathrm W}(R,g\cdot U)=S_{\mathrm W}(R,U)\) for every
\(g\in\mathcal G(R)\) and \(U\in\mathscr U(R)\).
\end{proof}

\begin{assumption}[Wilson comparison remainder]
\label{ass:wilson-comparison-remainder}
Assume there exist measurable functions
\[
c_{N,T}:\mathscr R_{N,T}\to\R,
\qquad
\Delta_{N,T}:\mathscr X_{N,T}\to\R,
\]
such that for \(\lambda_{N,T}\)-almost every \(R\) and \(\nu_R\)-almost every
\(U\),
\[
S_{\mathrm{eff}}^{(N,T)}(R,U)
=
S_{\mathrm W}(R,U)+\Delta_{N,T}(R,U)+c_{N,T}(R).
\]
Assume also that \(\Delta_{N,T}(R,\cdot)\) is gauge-invariant for
\(\lambda_{N,T}\)-almost every \(R\), and that its oscillation satisfies
\[
\operatorname{osc}_U \Delta_{N,T}(R,\cdot)
:=
\sup_{U\in\mathscr U(R)}\Delta_{N,T}(R,U)
-\inf_{U\in\mathscr U(R)}\Delta_{N,T}(R,U)
\le \varepsilon_{N,T}(R)
\]
for some measurable error profile \(\varepsilon_{N,T}(R)\ge 0\).
\end{assumption}

\begin{theorem}[Fiberwise comparison with the Wilson measure]
\label{thm:fiberwise-wilson-comparison}
Under Assumptions~\ref{ass:fiberwise-absolute-continuity} and
\ref{ass:wilson-comparison-remainder}, for \(\lambda_{N,T}\)-almost every
regulator \(R\), the induced fiber measure \(\mu_{N,T}^R\) and the Wilson fiber
measure \(\mu_W^R\) are mutually absolutely continuous, with density ratio
satisfying
\[
e^{-\varepsilon_{N,T}(R)}
\le
\frac{d\mu_{N,T}^R}{d\mu_W^R}(U)
\le
e^{\varepsilon_{N,T}(R)}
\qquad
\text{for }\mu_W^R\text{-almost every }U.
\]
Consequently,
\[
\|\mu_{N,T}^R-\mu_W^R\|_{\mathrm{TV}}
\le \sinh\!\bigl(\varepsilon_{N,T}(R)\bigr)
\le e^{\varepsilon_{N,T}(R)}-1.
\]
For every bounded measurable \(f\) on \(\mathscr U(R)\),
\[
\left|
\int f\,d\mu_{N,T}^R-\int f\,d\mu_W^R
\right|
\le
2\|f\|_\infty\,\sinh\!\bigl(\varepsilon_{N,T}(R)\bigr).
\]
\end{theorem}

\begin{proof}
Fix \(R\) in the full-measure set where the assumptions hold. By
Theorem~\ref{thm:exact-effective-action} and
Definition~\ref{def:fiberwise-wilson-measure},
\[
\frac{d\mu_{N,T}^R}{d\mu_W^R}(U)
=
\frac{Z_W(R)}{e^{c_{N,T}(R)}}e^{-\Delta_{N,T}(R,U)}.
\]
Let
\[
m(R):=\inf_{U}\Delta_{N,T}(R,U),\qquad
M(R):=\sup_{U}\Delta_{N,T}(R,U).
\]
Then \(M(R)-m(R)\le \varepsilon_{N,T}(R)\) and
\[
e^{-M(R)}
\le
\int e^{-\Delta_{N,T}(R,U)}\,\mu_W^R(dU)
\le
e^{-m(R)}.
\]
Since
\[
\frac{Z_W(R)}{e^{c_{N,T}(R)}}=
\left(
\int e^{-\Delta_{N,T}(R,U)}\,\mu_W^R(dU)
\right)^{-1},
\]
we obtain
\[
e^{m(R)-M(R)}
\le
\frac{d\mu_{N,T}^R}{d\mu_W^R}(U)
\le
e^{M(R)-m(R)},
\]
hence the displayed density bound. The total-variation estimate then follows
from
\[
\|\mu-\nu\|_{\mathrm{TV}}
=
\frac12\int\left|\frac{d\mu}{d\nu}-1\right|\,d\nu,
\]
and the observable estimate is the standard inequality
\[
\left|\int f\,d\mu-\int f\,d\nu\right|
\le 2\|f\|_\infty\|\mu-\nu\|_{\mathrm{TV}}.
\]
\end{proof}

\begin{corollary}[Mixture-level Wilson comparison]
\label{cor:mixture-level-wilson-comparison}
Under Assumptions~\ref{ass:fiberwise-absolute-continuity} and
\ref{ass:wilson-comparison-remainder},
\[
\|\mu_{N,T}-\mu_W^{(N,T)}\|_{\mathrm{TV}}
\le
\int_{\mathscr R_{N,T}}
\sinh\!\bigl(\varepsilon_{N,T}(R)\bigr)\,\lambda_{N,T}(dR).
\]
Hence for every bounded measurable gauge observable \(\Phi\) on
\(\mathscr X_{N,T}\),
\[
\left|
\int\Phi\,d\mu_{N,T}-\int\Phi\,d\mu_W^{(N,T)}
\right|
\le
2\|\Phi\|_\infty
\int_{\mathscr R_{N,T}}
\sinh\!\bigl(\varepsilon_{N,T}(R)\bigr)\,\lambda_{N,T}(dR).
\]
\end{corollary}

\begin{proof}
Disintegrate both laws over the common regulator law \(\lambda_{N,T}\) and
integrate the fiberwise total-variation estimate from
Theorem~\ref{thm:fiberwise-wilson-comparison}.
\end{proof}

\begin{corollary}[Asymptotic Wilson equivalence criterion]
\label{cor:asymptotic-wilson-equivalence}
Let \((N_k,T_k)\) be a sequence of finite regulators. If
\[
\int_{\mathscr R_{N_k,T_k}}
\sinh\!\bigl(\varepsilon_{N_k,T_k}(R)\bigr)\,\lambda_{N_k,T_k}(dR)
\longrightarrow 0,
\]
then
\[
\|\mu_{N_k,T_k}-\mu_W^{(N_k,T_k)}\|_{\mathrm{TV}}\to 0.
\]
In particular, every bounded gauge observable has the same limit under the
induced law and the mixed Wilson law.
\end{corollary}

\begin{proof}
Immediate from Corollary~\ref{cor:mixture-level-wilson-comparison}.
\end{proof}

\begin{remark}[What Stage I now proves]
\label{rem:what-stage1-proves}
After the package above, the unconditional finite-volume theorem package is:
\begin{enumerate}[leftmargin=1.5em,label=\textnormal{(\roman*)}]
\item the realized-regulator space is standard Borel;
\item the finite gauge fibers are compact Polish spaces;
\item the induced gauge law exists on the total gauge space;
\item the induced gauge law disintegrates over realized regulators; and
\item Wilson loops and the Wilson action are bounded measurable observables.
\end{enumerate}
The conditional finite-volume Yang--Mills-type package is:
\begin{enumerate}[leftmargin=1.5em,label=\textnormal{(\roman*)}]
\item under Assumption~\ref{ass:gauge-equivariant-lift}, the conditional fiber
laws are gauge-invariant;
\item under Assumption~\ref{ass:fiberwise-absolute-continuity}, the exact
effective action exists;
\item under Assumption~\ref{ass:wilson-comparison-remainder}, the induced fiber
law is exponentially close to the Wilson fiber law in total variation; and
\item if the remainder vanishes under refinement, the induced finite gauge
theory becomes asymptotically equivalent to a mixed Wilson gauge theory.
\end{enumerate}
\end{remark}

\begin{remark}[Exact insertion plan]
\label{rem:stage1-insertion-plan}
The old single theorem asserting only existence of a pushforward gauge law has
been replaced by the staged package:
\begin{enumerate}[leftmargin=1.5em,label=\textnormal{(\roman*)}]
\item realized regulator-data space;
\item fiber gauge space;
\item total gauge space;
\item induced law plus disintegration;
\item fiber Haar measure;
\item fiberwise absolute continuity;
\item exact effective action;
\item fiber Wilson measure;
\item Wilson comparison remainder; and
\item fiberwise and mixed-law Wilson comparison, including asymptotic equivalence.
\end{enumerate}
The Wilson-action formulas above are now read as defining the reference Wilson
fiber measures \(\mu_W^R\), not as the final finite gauge law of the theory.
The continuum comparison theorem remains later in the paper and plays the
strictly secondary role of comparing local small-loop expectations to the
classical continuum field.
\end{remark}

\begin{remark}[Remaining model-side hypotheses]
\label{rem:remaining-model-side-hypotheses}
Stage I is rigorous as stated, but two assumptions remain to be discharged from
the Fractal-Gas dynamics.
\begin{enumerate}[leftmargin=1.5em,label=\textnormal{(\roman*)}]
\item To prove Assumption~\ref{ass:fiberwise-absolute-continuity}, one needs a
model-specific argument that conditional on a realized regulator, the induced
edge variables admit a density relative to product Haar measure. Natural routes
include explicit edge noise in Lie-algebra coordinates before exponentiation,
smooth nondegeneracy of the edge phase/color map, or Malliavin-type arguments
for smooth nondegenerate Gaussian driving noise.
\item To prove Assumption~\ref{ass:wilson-comparison-remainder}, one must
identify the exact effective action and show that it is close to the Wilson
action by expanding the conditional log-density, isolating the plaquette-local
quadratic contribution, identifying it with the Wilson action, and proving that
the remainder is gauge-invariant with small oscillation. A sufficient route is
to show that, conditional on the realized regulator, the edge law is a
quasi-local Gibbs measure whose leading plaquette term is the Wilson action and
whose higher interactions are uniformly small.
\end{enumerate}
\end{remark}

\begin{remark}[Why Stage I is the correct target]
\label{rem:why-stage1-correct}
Before Stage I, the paper had a random regulator, random gauge variables,
Wilson loops, and a pushforward law. After Stage I, it has a finite-volume
gauge measure theory on random regulators, exact conditional gauge measures, an
exact effective action under a clean absolute-continuity hypothesis, and a
rigorous comparison theorem showing when the swarm-induced law is Wilson-like.
That is the correct launching point for the later Euclidean-QFT and mass-gap
stages.
\end{remark}

\begin{remark}[Final Stage I summary]
\label{rem:stage1-final-summary}
The finite random gauge construction no longer stops at the statement that a
pushforward gauge law exists. It now proves the stronger structural statement
that the realized Fractal regulator supports a finite-volume gauge measure
theory, disintegrated over realized regulators, with an exact effective action
and a rigorous criterion for Wilson equivalence.
\end{remark}

\section{Classical continuum comparison}

The stationary mean-field field \(A^\infty\), \(F^\infty\) from Section~3
determines the continuum Yang--Mills density appearing in the local theorem
below.

\begin{definition}[Localized continuum Yang--Mills functional]
\label{def:ym-action}
For \(\chi\in C_c(\R^4)\), define
\[
S_{\mathrm{YM},\chi}^{\mathrm{cont}}[F^\infty]
:=
\sum_{0\le\mu<\nu\le3}\int_{\R^4}
\chi(x)\,\tr\!\bigl((F_{\mu\nu}^\infty(x))^2\bigr)\,\dd x.
\]
Equivalently,
\[
S_{\mathrm{YM},\chi}^{\mathrm{cont}}[F^\infty]
=
\frac14\int_{\R^4}\chi(x)\sum_{\mu,\nu=0}^3\sum_{a=1}^8
F_{\mu\nu}^{\infty,a}(x)F_{\mu\nu}^{\infty,a}(x)\,\dd x.
\]
\end{definition}

\begin{assumption}[Local small-loop comparison regime]
\label{ass:local-comparison-regime}
Fix a compact window \(\mcK\subset\R^4\). For each \(N,T\), suppose there is a
refinement scale \(r_{N,T,\mcK}\downarrow0\) and measurable orientation labels
\[
\mathrm{or}(P)=(\mu(P),\nu(P))\in\{(\mu,\nu):0\le\mu<\nu\le3\}
\]
for plaquettes with \(x_P(\omega)\in\mcK\), such that the following hold.
\begin{enumerate}[leftmargin=1.5em,label=\textnormal{(\roman*)}]
\item \textbf{Principal-log/small-field regime.}
For every \(\delta>0\) small enough,
\[
\mathbb P_{N,T}\!\left(
\sup_{\substack{P\in\mcP_{\mathrm{FG}}^{(N,T)}\\ x_P\in\mcK}}
\|g_3A_PF_P\|_{\mathrm{op}}\le\delta
\right)\longrightarrow1.
\]
\item \textbf{Weighted curvature identification in expectation.}
For every \(0\le\mu<\nu\le3\),
\[
\E\!\left[
\sum_{\substack{P\in\mcP_{\mathrm{FG}}^{(N,T)}\\
x_P\in\mcK,\ \mathrm{or}(P)=(\mu,\nu)}}
A_P^2
\left|
\tr\!\bigl((F_P)^2\bigr)
-\tr\!\bigl((F_{\mu\nu}^\infty(x_P))^2\bigr)
\right|
\right]
\longrightarrow0.
\]
\item \textbf{Orientation-wise Riemann convergence in expectation.}
For every \(\varphi\in C_c(\R^4)\) supported in \(\mcK\) and every
\(0\le\mu<\nu\le3\),
\[
\E\!\left[
\sum_{\substack{P\in\mcP_{\mathrm{FG}}^{(N,T)}\\ \mathrm{or}(P)=(\mu,\nu)}}
A_P^2\,\varphi(x_P)
\right]
\longrightarrow
\int_{\R^4}\varphi(x)\,\dd x.
\]
\end{enumerate}
\end{assumption}

\begin{theorem}[Localized classical comparison of Wilson observables]
\label{thm:continuum-wilson}
Assume Assumption~\ref{ass:local-comparison-regime}. For \(\chi\in C_c(\R^4)\),
define the localized random Wilson functional
\[
S_{\mathrm W,\chi}^{(N,T)}(\omega)
:=
\frac{6}{g_3^2}
\sum_{P\in\mcP_{\mathrm{FG}}^{(N,T)}(\omega)}
\chi(x_P(\omega))
\left(
1-\frac13\Re\tr U_P(\omega)
\right).
\]
Then
\[
\E_{\mathbb P_{N,T}}\bigl[S_{\mathrm W,\chi}^{(N,T)}\bigr]
\longrightarrow
S_{\mathrm{YM},\chi}^{\mathrm{cont}}[F^\infty].
\]
\end{theorem}

\begin{proof}
Let \(\mcK:=\supp\chi\). By the small-field part of
Assumption~\ref{ass:local-comparison-regime} and
Proposition~\ref{prop:random-plaquette-wilson-expansion}, the plaquette term
admits the expansion
\[
1-\frac13\Re\tr U_P
=
\frac{g_3^2A_P^2}{6}\tr\!\bigl((F_P)^2\bigr)+R_P,
\]
with a remainder whose expectation over plaquettes in \(\mcK\) vanishes as the
refinement scale tends to \(0\). Therefore
\begin{align*}
\E\bigl[S_{\mathrm W,\chi}^{(N,T)}\bigr]
&=
\E\!\left[
\sum_{P\in\mcP_{\mathrm{FG}}^{(N,T)}}
\chi(x_P)A_P^2\tr\!\bigl((F_P)^2\bigr)
\right]
+o(1).
\end{align*}
Split the last sum by orientation. For each \(0\le\mu<\nu\le3\),
Assumption~\ref{ass:local-comparison-regime}\textnormal{(ii)} yields
\begin{align*}
&\E\!\left[
\sum_{\substack{P\in\mcP_{\mathrm{FG}}^{(N,T)}\\ \mathrm{or}(P)=(\mu,\nu)}}
\chi(x_P)A_P^2\tr\!\bigl((F_P)^2\bigr)
\right]\\
&\qquad=
\E\!\left[
\sum_{\substack{P\in\mcP_{\mathrm{FG}}^{(N,T)}\\ \mathrm{or}(P)=(\mu,\nu)}}
\chi(x_P)A_P^2\tr\!\bigl((F_{\mu\nu}^\infty(x_P))^2\bigr)
\right]
+o(1).
\end{align*}
Applying Assumption~\ref{ass:local-comparison-regime}\textnormal{(iii)} to
\[
\varphi_{\mu\nu}(x):=\chi(x)\tr\!\bigl((F_{\mu\nu}^\infty(x))^2\bigr)
\]
gives
\[
\E\!\left[
\sum_{\substack{P\in\mcP_{\mathrm{FG}}^{(N,T)}\\ \mathrm{or}(P)=(\mu,\nu)}}
\chi(x_P)A_P^2\tr\!\bigl((F_{\mu\nu}^\infty(x_P))^2\bigr)
\right]
\longrightarrow
\int_{\R^4}\chi(x)\tr\!\bigl((F_{\mu\nu}^\infty(x))^2\bigr)\,\dd x.
\]
Summing over the six unordered pairs \(\mu<\nu\) gives the claim.
\end{proof}

\section{Exported theorem package for downstream papers}

\begin{theorem}[Export package of the revised Paper 03b]
\label{thm:paper2-export}
The theorem package exported to later papers in the series is the following.
\begin{enumerate}[leftmargin=1.5em,label=\textnormal{(\roman*)}]
\item \textbf{Stochastic substrate from \texttt{01\_convergence}.}
Section~2 supplies the trajectory laws, stationary limits, and compactness
package needed to speak about finite swarm histories and their thermodynamic
limits.
\item \textbf{Realized Fractal regulator.}
Definition~\ref{def:history-path-law},
Definition~\ref{def:realized-fractal-regulator},
Definition~\ref{def:fractal-plaquette-geometry}, and
Proposition~\ref{prop:fractal-regulator-measurable} construct the finite random
oriented \(2\)-complex extracted from each realized run.
\item \textbf{Fibered finite-volume gauge measure theory.}
Definition~\ref{def:realized-node-color},
Definition~\ref{def:realized-edge-transports},
Definition~\ref{def:realized-wilson-observables},
Definition~\ref{def:realized-regulator-data-space},
Definition~\ref{def:fiber-gauge-configuration-space},
Definition~\ref{def:total-fibered-gauge-space},
Theorem~\ref{thm:induced-finite-volume-gauge-law}, and
Theorem~\ref{thm:disintegration-over-regulators} construct the realized edge
transports, plaquette holonomies, Wilson observables, the total finite gauge
space \(\mathscr X_{N,T}\), the induced law
\[
\mu_{N,T},
\]
and its disintegration over realized regulators.
\item \textbf{Exact finite-volume gauge action and Wilson comparison.}
Definition~\ref{def:fiberwise-reference-haar},
Assumption~\ref{ass:fiberwise-absolute-continuity},
Theorem~\ref{thm:exact-effective-action},
Definition~\ref{def:fiberwise-wilson-measure},
Assumption~\ref{ass:wilson-comparison-remainder},
Theorem~\ref{thm:fiberwise-wilson-comparison}, and
Corollary~\ref{cor:asymptotic-wilson-equivalence} identify the exact
effective action whenever the conditional laws are absolutely continuous and
state the rigorous criterion under which the induced finite gauge theory is
Wilson-like.
\item \textbf{Continuum comparison field.}
Definition~\ref{def:continuum-color-field},
Definition~\ref{def:stationary-mf-color},
Proposition~\ref{prop:stationary-mf-color-smooth},
Proposition~\ref{prop:continuum-connection-regularity}, and
Proposition~\ref{prop:continuum-energy-tail} construct the smooth stationary
comparison connection \(A^\infty\) and curvature \(F^\infty\).
\item \textbf{Classical comparison theorem in expectation.}
Definition~\ref{def:ym-action},
Assumption~\ref{ass:local-comparison-regime}, and
Theorem~\ref{thm:continuum-wilson} identify the regime in which expected
localized Wilson observables recover the continuum Yang--Mills density.
\item \textbf{04b interface.}
The downstream Euclidean paper should import the random regulator law
\(\lambda_{N,T}\), the total gauge law \(\mu_{N,T}\), its disintegration
\(\{\mu_{N,T}^R\}\), the reference Wilson fibers \(\mu_W^R\), and the Wilson
comparison package directly, together with the local continuum-density theorem
proved here.
\end{enumerate}
\end{theorem}

\begin{proof}
Each item is exactly the labeled construction or theorem cited in its
statement. The final item is the interface summary extracted from the preceding
five items.
\end{proof}

\begin{remark}[How later papers use this export package]
Paper~04b should import the regulator law \(\lambda_{N,T}\), the induced total
gauge law \(\mu_{N,T}\), the conditional fiber laws \(\mu_{N,T}^R\), the
reference Wilson fibers \(\mu_W^R\), and the corresponding Wilson and plaquette
observables from Theorem~\ref{thm:paper2-export}.
\end{remark}

\begin{thebibliography}{9}

\bibitem{billingsley1999convergence}
P.~Billingsley,
\emph{Convergence of Probability Measures},
2nd ed.,
Wiley, New York, 1999.

\bibitem{diaconis1980finite}
P.~Diaconis and D.~Freedman,
``Finite exchangeable sequences,''
\emph{The Annals of Probability} \textbf{8} (1980), no.~4, 745--764.

\end{thebibliography}


\end{document}
